\section{\OmniASR Language Coverage}
\label{sec:appending}

{
    \tiny\renewcommand{\arraystretch}{.8}
    \begin{adjustbox}{max width=\textwidth}
    \begin{tabular}{llc}
        \toprule
        \textbf{Code} & \textbf{Name} & \textbf{Res}\\
        \midrule
        \texttt{aae\_Latn} & Arbëreshë Albanian & L \\
        \texttt{aal\_Latn} & Afade & L \\
        \texttt{abb\_Latn} & Bankon & L \\
        \texttt{abi\_Latn} & Abidji & M \\
        \texttt{abk\_Cyrl} & Abkhazian & M \\
        \texttt{abn\_Latn} & Abua & L \\
        \texttt{abp\_Latn} & Abellen Ayta & M \\
        \texttt{abr\_Latn} & Abron & L \\
        \texttt{abs\_Latn} & Ambonese Malay & L \\
        \texttt{aca\_Latn} & Achagua & M \\
        \texttt{acd\_Latn} & Gikyode & M \\
        \texttt{ace\_Latn} & Achinese & M \\
        \texttt{acf\_Latn} & Saint Lucian Creole French & M \\
        \texttt{ach\_Latn} & Acoli & M \\
        \texttt{acm\_Arab} & Mesopotamian Arabic & L \\
        \texttt{acn\_Latn} & Achang & M \\
        \texttt{acr\_Latn} & Achi & M \\
        \texttt{acu\_Latn} & Achuar-Shiwiar & M \\
        \texttt{acw\_Arab} & Hijazi Arabic & M \\
        \texttt{ade\_Latn} & Adele & M \\
        \texttt{adh\_Latn} & Adhola & M \\
        \texttt{adj\_Latn} & Adioukrou & M \\
        \texttt{adx\_Tibt} & Amdo Tibetan & H \\
        \texttt{ady\_Cyrl} & Adyghe & M \\
        \texttt{aeb\_Arab} & Tunisian Arabic & M \\
        \texttt{aec\_Arab} & Saidi Arabic & L \\
        \texttt{aeu\_Latn} & Akeu & M \\
        \texttt{afb\_Arab} & Gulf Arabic & M \\
        \texttt{afo\_Latn} & Eloyi & L \\
        \texttt{afr\_Latn} & Afrikaans & H \\
        \texttt{agd\_Latn} & Agarabi & M \\
        \texttt{agg\_Latn} & Angor & L \\
        \texttt{agn\_Latn} & Agutaynen & M \\
        \texttt{agr\_Latn} & Aguaruna & M \\
        \texttt{agu\_Latn} & Aguacateco & M \\
        \texttt{agx\_Cyrl} & Aghul & L \\
        \texttt{aha\_Latn} & Ahanta & M \\
        \texttt{ahk\_Latn} & Akha & M \\
        \texttt{ahl\_Latn} & Igo & L \\
        \texttt{ahs\_Latn} & Ashe & L \\
        \texttt{aia\_Latn} & Arosi & M \\
        \texttt{ajg\_Latn} & Aja (Benin) & L \\
        \texttt{aka\_Latn} & Akan & M \\
        \texttt{akb\_Latn} & Batak Angkola & M \\
        \texttt{ake\_Latn} & Akawaio & M \\
        \texttt{akp\_Latn} & Siwu & M \\
        \texttt{ala\_Latn} & Alago & L \\
        \texttt{alj\_Latn} & Alangan & M \\
        \texttt{aln\_Latn} & Gheg Albanian & L \\
        \texttt{alo\_Latn} & Larike-Wakasihu & L \\
        \texttt{alp\_Latn} & Alune & M \\
        \texttt{als\_Latn} & Tosk Albanian & M \\
        \texttt{alt\_Cyrl} & Southern Altai & M \\
        \texttt{alz\_Latn} & Alur & M \\
        \texttt{ame\_Latn} & Yanesha' & H \\
        \texttt{amf\_Latn} & Hamer-Banna & M \\
        \texttt{amh\_Ethi} & Amharic & H \\
        \texttt{ami\_Latn} & Amis & H \\
        \texttt{amk\_Latn} & Ambai & H \\
        \texttt{amu\_Latn} & Guerrero Amuzgo & L \\
        \texttt{anc\_Latn} & Ngas & L \\
        \texttt{ank\_Latn} & Goemai & L \\
        \texttt{ann\_Latn} & Obolo & M \\
        \texttt{anp\_Deva} & Angika & L \\
        \texttt{anw\_Latn} & Anaang & L \\
        \texttt{any\_Latn} & Anyin & M \\
        \texttt{aom\_Latn} & Ömie & L \\
        \texttt{aoz\_Latn} & Uab Meto & L \\
        \texttt{apb\_Latn} & Sa'a & M \\
        \texttt{apc\_Arab} & Levantine Arabic & L \\
        \texttt{apd\_Arab} & Sudanese Arabic & L \\
        \texttt{apr\_Latn} & Arop-Lokep & M \\
        \texttt{arb\_Arab} & Standard Arabic & H \\
        \texttt{arg\_Latn} & Aragonese & M \\
        \texttt{arl\_Latn} & Arabela & H \\
        \texttt{arq\_Arab} & Algerian Arabic & L \\
        \texttt{ars\_Arab} & Najdi Arabic & M \\
        \texttt{ary\_Arab} & Moroccan Arabic & L \\
        \texttt{arz\_Arab} & Egyptian Arabic & L \\
        \texttt{asa\_Latn} & Asu (Tanzania) & M \\
        \texttt{asg\_Latn} & Cishingini & M \\
        \texttt{asm\_Beng} & Assamese & H \\
        \texttt{ast\_Latn} & Asturian & L \\
        \texttt{ata\_Latn} & Pele-Ata & M \\
        \texttt{atb\_Latn} & Zaiwa & M \\
        \texttt{atg\_Latn} & Ivbie North-Okpela-Arhe & M \\
        \texttt{ati\_Latn} & Attié & M \\
        \texttt{atq\_Latn} & Aralle-Tabulahan & H \\
        \texttt{ava\_Cyrl} & Avaric & M \\
        \texttt{avn\_Latn} & Avatime & M \\
        \texttt{avu\_Latn} & Avokaya & M \\
        \texttt{awa\_Deva} & Awadhi & M \\
        \texttt{awb\_Latn} & Awa (Papua New Guinea) & M \\
        \texttt{awo\_Latn} & Awak & L \\
        \texttt{ayl\_Arab} & Libyan Arabic & M \\
        \texttt{ayo\_Latn} & Ayoreo & H \\
        \texttt{ayp\_Arab} & North Mesopotamian Arabic & L \\
        \texttt{ayr\_Latn} & Central Aymara & M \\
        \texttt{ayz\_Latn} & Mai Brat & L \\
        \texttt{aze\_Arab} & Azerbaijani & M \\
        \texttt{aze\_Cyrl} & Azerbaijani & M \\
        \texttt{aze\_Latn} & Azerbaijani & M \\
        \texttt{azg\_Latn} & San Pedro Amuzgos Amuzgo & M \\
        \texttt{azz\_Latn} & Highland Puebla Nahuatl & M \\
        \texttt{bag\_Latn} & Tuki & L \\
        \texttt{bak\_Cyrl} & Bashkir & H \\
        \texttt{bam\_Latn} & Bambara & M \\
        \texttt{ban\_Latn} & Balinese & M \\
        \texttt{bao\_Latn} & Waimaha & M \\
        \texttt{bas\_Latn} & Basa (Cameroon) & L \\
        \texttt{bav\_Latn} & Vengo & M \\
        \texttt{bax\_Latn} & Bamun & L \\
        \texttt{bba\_Latn} & Baatonum & M \\
        \texttt{bbb\_Latn} & Barai & H \\
        \texttt{bbc\_Latn} & Batak Toba & M \\
        \texttt{bbj\_Latn} & Ghomálá' & L \\
        \texttt{bbl\_Geor} & Bats & L \\
        \texttt{bbo\_Latn} & Northern Bobo Madaré & M \\
        \texttt{bbu\_Latn} & Kulung (Nigeria) & L \\
        \texttt{bcc\_Arab} & Southern Balochi & M \\
        \texttt{bcc\_Latn} & Southern Balochi & M \\
        \texttt{bce\_Latn} & Bamenyam & L \\
        \texttt{bci\_Latn} & Baoulé & L \\
        \texttt{bcl\_Latn} & Central Bikol & M \\
        \texttt{bcs\_Latn} & Kohumono & L \\
        \bottomrule
    \end{tabular}
    \quad
    \begin{tabular}{llc}
        \toprule
        \textbf{Code} & \textbf{Name} & \textbf{Res}\\
        \midrule
        \texttt{bcw\_Latn} & Bana & M \\
        \texttt{bcy\_Latn} & Bacama & L \\
        \texttt{bcz\_Latn} & Bainouk-Gunyaamolo & L \\
        \texttt{bda\_Latn} & Bayot & L \\
        \texttt{bde\_Latn} & Bade & L \\
        \texttt{bdg\_Latn} & Bonggi & M \\
        \texttt{bdh\_Latn} & Baka (South Sudan) & H \\
        \texttt{bdm\_Latn} & Buduma & L \\
        \texttt{bdq\_Latn} & Bahnar & M \\
        \texttt{bdu\_Latn} & Oroko & M \\
        \texttt{beb\_Latn} & Bebele & L \\
        \texttt{beh\_Latn} & Biali & M \\
        \texttt{bel\_Cyrl} & Belarusian & H \\
        \texttt{bem\_Latn} & Bemba (Zambia) & M \\
        \texttt{ben\_Beng} & Bengali & H \\
        \texttt{bep\_Latn} & Besoa & M \\
        \texttt{bew\_Latn} & Betawi & M \\
        \texttt{bex\_Latn} & Jur Modo & M \\
        \texttt{bfa\_Latn} & Bari & M \\
        \texttt{bfd\_Latn} & Bafut & L \\
        \texttt{bfo\_Latn} & Malba Birifor & M \\
        \texttt{bft\_Arab} & Balti & M \\
        \texttt{bfy\_Deva} & Bagheli & M \\
        \texttt{bfz\_Deva} & Mahasu Pahari & M \\
        \texttt{bgc\_Deva} & Haryanvi & M \\
        \texttt{bgp\_Arab} & Eastern Balochi & L \\
        \texttt{bgq\_Deva} & Bagri & M \\
        \texttt{bgr\_Latn} & Bawm Chin & M \\
        \texttt{bgt\_Latn} & Bughotu & M \\
        \texttt{bgw\_Deva} & Bhatri & M \\
        \texttt{bha\_Deva} & Bharia & L \\
        \texttt{bhb\_Deva} & Bhili & L \\
        \texttt{bhh\_Cyrl} & Bukharic & L \\
        \texttt{bho\_Deva} & Bhojpuri & L \\
        \texttt{bhp\_Latn} & Bima & L \\
        \texttt{bht\_Deva} & Bhattiyali & M \\
        \texttt{bhz\_Latn} & Bada (Indonesia) & M \\
        \texttt{bib\_Latn} & Bissa & M \\
        \texttt{bim\_Latn} & Bimoba & M \\
        \texttt{bis\_Latn} & Bislama & M \\
        \texttt{biv\_Latn} & Southern Birifor & M \\
        \texttt{bjj\_Deva} & Kanauji & L \\
        \texttt{bjk\_Latn} & Barok & L \\
        \texttt{bjn\_Latn} & Banjar & L \\
        \texttt{bjr\_Latn} & Binumarien & H \\
        \texttt{bjt\_Latn} & Balanta-Ganja & L \\
        \texttt{bjv\_Latn} & Bedjond & M \\
        \texttt{bjw\_Latn} & Bakwé & M \\
        \texttt{bjz\_Latn} & Baruga & M \\
        \texttt{bkd\_Latn} & Binukid & M \\
        \texttt{bkh\_Latn} & Bakoko & L \\
        \texttt{bkm\_Latn} & Kom (Cameroon) & L \\
        \texttt{bkv\_Latn} & Bekwarra & M \\
        \texttt{bky\_Latn} & Bokyi & L \\
        \texttt{ble\_Latn} & Balanta-Kentohe & L \\
        \texttt{blh\_Latn} & Kuwaa & M \\
        \texttt{blt\_Latn} & Tai Dam & M \\
        \texttt{blx\_Latn} & Mag-Indi Ayta & M \\
        \texttt{blz\_Latn} & Balantak & M \\
        \texttt{bmm\_Latn} & Northern Betsimisaraka Malagasy & M \\
        \texttt{bmq\_Latn} & Bomu & M \\
        \texttt{bmr\_Latn} & Muinane & H \\
        \texttt{bmu\_Latn} & Somba-Siawari & M \\
        \texttt{bmv\_Latn} & Bum & M \\
        \texttt{bng\_Beng} & Benga & M \\
        \texttt{bnm\_Latn} & Batanga & M \\
        \texttt{bnn\_Latn} & Bunun & L \\
        \texttt{bno\_Latn} & Bantoanon & H \\
        \texttt{bnp\_Latn} & Bola & M \\
        \texttt{bns\_Deva} & Bundeli & L \\
        \texttt{boa\_Latn} & Bora & M \\
        \texttt{bod\_Tibt} & Tibetan & H \\
        \texttt{boj\_Latn} & Anjam & H \\
        \texttt{bom\_Latn} & Berom & M \\
        \texttt{bor\_Latn} & Borôro & H \\
        \texttt{bos\_Latn} & Bosnian & H \\
        \texttt{bou\_Latn} & Bondei & L \\
        \texttt{bov\_Latn} & Tuwuli & M \\
        \texttt{box\_Latn} & Buamu & M \\
        \texttt{bpr\_Latn} & Koronadal Blaan & M \\
        \texttt{bps\_Latn} & Sarangani Blaan & M \\
        \texttt{bqc\_Latn} & Boko (Benin) & M \\
        \texttt{bqg\_Latn} & Bago-Kusuntu & L \\
        \texttt{bqi\_Arab} & Bakhtiari & L \\
        \texttt{bqj\_Latn} & Bandial & M \\
        \texttt{bqp\_Latn} & Busa & M \\
        \texttt{bra\_Deva} & Braj & L \\
        \texttt{bre\_Latn} & Breton & L \\
        \texttt{brh\_Arab} & Brahui & M \\
        \texttt{bri\_Latn} & Mokpwe & L \\
        \texttt{bru\_Latn} & Eastern Bru & M \\
        \texttt{brx\_Deva} & Bodo (India) & L \\
        \texttt{bsc\_Latn} & Bassari & M \\
        \texttt{bsh\_Arab} & Kati & L \\
        \texttt{bsj\_Latn} & Bangwinji & L \\
        \texttt{bsk\_Latn} & Burushaski & L \\
        \texttt{bsq\_Latn} & Bassa & M \\
        \texttt{bss\_Latn} & Akoose & M \\
        \texttt{bsy\_Latn} & Sabah Bisaya & L \\
        \texttt{btd\_Latn} & Batak Dairi & M \\
        \texttt{btm\_Latn} & Batak Mandailing & L \\
        \texttt{bts\_Latn} & Batak Simalungun & M \\
        \texttt{btt\_Latn} & Bete-Bendi & M \\
        \texttt{btv\_Arab} & Bateri & L \\
        \texttt{btx\_Latn} & Batak Karo & M \\
        \texttt{bud\_Latn} & Ntcham & M \\
        \texttt{bug\_Latn} & Buginese & L \\
        \texttt{bul\_Cyrl} & Bulgarian & H \\
        \texttt{bum\_Latn} & Bulu (Cameroon) & L \\
        \texttt{buo\_Latn} & Terei & L \\
        \texttt{bus\_Latn} & Bokobaru & M \\
        \texttt{bux\_Latn} & Boghom & L \\
        \texttt{bvb\_Latn} & Bube & L \\
        \texttt{bvc\_Latn} & Baelelea & M \\
        \texttt{bvz\_Latn} & Bauzi & H \\
        \texttt{bwq\_Latn} & Southern Bobo Madaré & M \\
        \texttt{bwr\_Latn} & Bura-Pabir & L \\
        \texttt{bwu\_Latn} & Buli (Ghana) & M \\
        \texttt{bxf\_Latn} & Bilur & L \\
        \texttt{bxk\_Latn} & Bukusu & L \\
        \texttt{byc\_Latn} & Ubaghara & L \\
        \texttt{byr\_Latn} & Baruya & H \\
        \texttt{bys\_Latn} & Burak & L \\
        \texttt{byv\_Latn} & Medumba & M \\
        \texttt{byx\_Latn} & Qaqet & L \\
        \bottomrule
    \end{tabular}
    \quad
    \begin{tabular}{llc}
        \toprule
        \textbf{Code} & \textbf{Name} & \textbf{Res}\\
        \midrule
        \texttt{bzh\_Latn} & Mapos Buang & M \\
        \texttt{bzi\_Thai} & Bisu & M \\
        \texttt{bzj\_Latn} & Belize Kriol English & M \\
        \texttt{bzw\_Latn} & Basa (Nigeria) & L \\
        \texttt{caa\_Latn} & Chortí & M \\
        \texttt{cab\_Latn} & Garifuna & M \\
        \texttt{cac\_Latn} & Chuj & H \\
        \texttt{cak\_Latn} & Kaqchikel & H \\
        \texttt{cap\_Latn} & Chipaya & M \\
        \texttt{car\_Latn} & Galibi Carib & M \\
        \texttt{cas\_Latn} & Tsimané & M \\
        \texttt{cat\_Latn} & Catalan & H \\
        \texttt{cax\_Latn} & Chiquitano & M \\
        \texttt{cbc\_Latn} & Carapana & M \\
        \texttt{cbi\_Latn} & Chachi & H \\
        \texttt{cbr\_Latn} & Cashibo-Cacataibo & H \\
        \texttt{cbs\_Latn} & Cashinahua & M \\
        \texttt{cbt\_Latn} & Chayahuita & M \\
        \texttt{cbu\_Latn} & Candoshi-Shapra & M \\
        \texttt{cbv\_Latn} & Cacua & M \\
        \texttt{cce\_Latn} & Chopi & M \\
        \texttt{ccg\_Latn} & Samba Daka & L \\
        \texttt{cco\_Latn} & Comaltepec Chinantec & M \\
        \texttt{cdj\_Deva} & Churahi & M \\
        \texttt{cdo\_Hans} & Min Dong Chinese & L \\
        \texttt{ceb\_Latn} & Cebuano & H \\
        \texttt{ceg\_Latn} & Chamacoco & M \\
        \texttt{cek\_Latn} & Eastern Khumi Chin & H \\
        \texttt{cen\_Latn} & Cen & L \\
        \texttt{ces\_Latn} & Czech & H \\
        \texttt{cfa\_Latn} & Dijim-Bwilim & L \\
        \texttt{cfm\_Latn} & Falam Chin & M \\
        \texttt{cgc\_Latn} & Kagayanen & M \\
        \texttt{cgg\_Latn} & Chiga & M \\
        \texttt{che\_Cyrl} & Chechen & M \\
        \texttt{chf\_Latn} & Tabasco Chontal & L \\
        \texttt{chq\_Latn} & Quiotepec Chinantec & L \\
        \texttt{chv\_Cyrl} & Chuvash & M \\
        \texttt{chz\_Latn} & Ozumacín Chinantec & M \\
        \texttt{cjk\_Latn} & Chokwe & L \\
        \texttt{cjo\_Latn} & Ashéninka Pajonal & H \\
        \texttt{cjp\_Latn} & Cabécar & M \\
        \texttt{cjs\_Cyrl} & Shor & L \\
        \texttt{ckb\_Arab} & Central Kurdish & L \\
        \texttt{ckl\_Latn} & Cibak & L \\
        \texttt{cko\_Latn} & Anufo & M \\
        \texttt{ckr\_Latn} & Kairak & L \\
        \texttt{ckt\_Cyrl} & Chukot & L \\
        \texttt{cky\_Latn} & Cakfem-Mushere & L \\
        \texttt{cla\_Latn} & Ron & M \\
        \texttt{cle\_Latn} & Lealao Chinantec & M \\
        \texttt{cly\_Latn} & Eastern Highland Chatino & M \\
        \texttt{cme\_Latn} & Cerma & M \\
        \texttt{cmn\_Hans} & Mandarin Chinese & H \\
        \texttt{cmn\_Hant} & Mandarin Chinese & M \\
        \texttt{cmo\_Khmr} & Central Mnong & M \\
        \texttt{cmo\_Latn} & Central Mnong & M \\
        \texttt{cmr\_Latn} & Mro-Khimi Chin & M \\
        \texttt{cnh\_Latn} & Hakha Chin & M \\
        \texttt{cni\_Latn} & Asháninka & M \\
        \texttt{cnl\_Latn} & Lalana Chinantec & M \\
        \texttt{cnt\_Latn} & Tepetotutla Chinantec & M \\
        \texttt{coe\_Latn} & Koreguaje & M \\
        \texttt{cof\_Latn} & Colorado & H \\
        \texttt{cok\_Latn} & Santa Teresa Cora & H \\
        \texttt{con\_Latn} & Cofán & M \\
        \texttt{cor\_Latn} & Cornish & L \\
        \texttt{cot\_Latn} & Caquinte & H \\
        \texttt{cou\_Latn} & Wamey & M \\
        \texttt{cpa\_Latn} & Palantla Chinantec & M \\
        \texttt{cpb\_Latn} & Ucayali-Yurúa Ashéninka & H \\
        \texttt{cpu\_Latn} & Pichis Ashéninka & H \\
        \texttt{cpx\_Hans} & Pu-Xian Chinese & L \\
        \texttt{cpy\_Latn} & South Ucayali Ashéninka & L \\
        \texttt{crh\_Cyrl} & Crimean Tatar & M \\
        \texttt{crk\_Cans} & Plains Cree & M \\
        \texttt{crk\_Latn} & Plains Cree & M \\
        \texttt{crn\_Latn} & El Nayar Cora & M \\
        \texttt{crq\_Latn} & Iyo'wujwa Chorote & H \\
        \texttt{crs\_Latn} & Seselwa Creole French & M \\
        \texttt{crt\_Latn} & Iyojwa'ja Chorote & H \\
        \texttt{csk\_Latn} & Jola-Kasa & M \\
        \texttt{cso\_Latn} & Sochiapam Chinantec & M \\
        \texttt{ctd\_Latn} & Tedim Chin & M \\
        \texttt{cte\_Latn} & Tepinapa Chinantec & L \\
        \texttt{ctg\_Beng} & Chittagonian & M \\
        \texttt{ctl\_Latn} & Tlacoatzintepec Chinantec & L \\
        \texttt{cto\_Latn} & Emberá-Catío & L \\
        \texttt{ctu\_Latn} & Chol & L \\
        \texttt{cuc\_Latn} & Usila Chinantec & M \\
        \texttt{cui\_Latn} & Cuiba & H \\
        \texttt{cuk\_Latn} & San Blas Kuna & M \\
        \texttt{cul\_Latn} & Culina & H \\
        \texttt{cut\_Latn} & Teutila Cuicatec & L \\
        \texttt{cux\_Latn} & Tepeuxila Cuicatec & L \\
        \texttt{cwa\_Latn} & Kabwa & H \\
        \texttt{cwe\_Latn} & Kwere & M \\
        \texttt{cwt\_Latn} & Kuwaataay & M \\
        \texttt{cya\_Latn} & Nopala Chatino & M \\
        \texttt{cym\_Latn} & Welsh & H \\
        \texttt{daa\_Latn} & Dangaléat & M \\
        \texttt{dag\_Latn} & Dagbani & L \\
        \texttt{dah\_Latn} & Gwahatike & H \\
        \texttt{dan\_Latn} & Danish & H \\
        \texttt{dar\_Cyrl} & Dargwa & L \\
        \texttt{dav\_Latn} & Taita & L \\
        \texttt{dbd\_Latn} & Dadiya & L \\
        \texttt{dbj\_Latn} & Ida'an & L \\
        \texttt{dbq\_Latn} & Daba & H \\
        \texttt{dcc\_Arab} & Deccan & L \\
        \texttt{ddn\_Latn} & Dendi (Benin) & M \\
        \texttt{ded\_Latn} & Dedua & M \\
        \texttt{deg\_Latn} & Degema & L \\
        \texttt{des\_Latn} & Desano & M \\
        \texttt{deu\_Latn} & German & H \\
        \texttt{dga\_Latn} & Southern Dagaare & M \\
        \texttt{dgh\_Latn} & Dghwede & L \\
        \texttt{dgi\_Latn} & Northern Dagara & M \\
        \texttt{dgk\_Latn} & Dagba & M \\
        \texttt{dgo\_Deva} & Dogri (individual language) & M \\
        \texttt{dgr\_Latn} & Dogrib & M \\
        \texttt{dhi\_Deva} & Dhimal & M \\
        \texttt{did\_Latn} & Didinga & M \\
        \texttt{dig\_Latn} & Digo & M \\
        \texttt{dik\_Latn} & Southwestern Dinka & M \\
        \bottomrule
    \end{tabular}
    \end{adjustbox}
}

{
    \tiny\renewcommand{\arraystretch}{.8}
    \begin{adjustbox}{max width=\textwidth}
    \begin{tabular}{llc}
        \toprule
        \textbf{Code} & \textbf{Name} & \textbf{Res}\\
        \midrule
        \texttt{dip\_Latn} & Northeastern Dinka & M \\
        \texttt{div\_Thaa} & Dhivehi & L \\
        \texttt{dje\_Latn} & Zarma & M \\
        \texttt{djk\_Latn} & Eastern Maroon Creole & M \\
        \texttt{dmk\_Arab} & Domaaki & L \\
        \texttt{dml\_Arab} & Dameli & L \\
        \texttt{dnj\_Latn} & Dan & M \\
        \texttt{dnt\_Latn} & Mid Grand Valley Dani & H \\
        \texttt{dnw\_Latn} & Western Dani & H \\
        \texttt{dop\_Latn} & Lukpa & M \\
        \texttt{dos\_Latn} & Dogosé & M \\
        \texttt{dru\_Latn} & Rukai & L \\
        \texttt{dsb\_Latn} & Lower Sorbian & L \\
        \texttt{dsh\_Latn} & Daasanach & H \\
        \texttt{dtp\_Latn} & Kadazan Dusun & H \\
        \texttt{dts\_Latn} & Toro So Dogon & M \\
        \texttt{dty\_Deva} & Dotyali & L \\
        \texttt{dua\_Latn} & Duala & L \\
        \texttt{dug\_Latn} & Duruma & M \\
        \texttt{dwr\_Latn} & Dawro & M \\
        \texttt{dyi\_Latn} & Djimini Senoufo & M \\
        \texttt{dyo\_Latn} & Jola-Fonyi & M \\
        \texttt{dyu\_Latn} & Dyula & H \\
        \texttt{dzg\_Latn} & Dazaga & L \\
        \texttt{dzo\_Tibt} & Dzongkha & M \\
        \texttt{ebu\_Latn} & Embu & L \\
        \texttt{ego\_Latn} & Eggon & L \\
        \texttt{eip\_Latn} & Eipomek & H \\
        \texttt{eiv\_Latn} & Askopan & L \\
        \texttt{eka\_Latn} & Ekajuk & M \\
        \texttt{ekk\_Latn} & Standard Estonian & H \\
        \texttt{eko\_Latn} & Koti & L \\
        \texttt{ekr\_Latn} & Yace & L \\
        \texttt{ell\_Grek} & Modern Greek & H \\
        \texttt{ell\_Grek\_cypr1249} & Cypriot Greek & L \\
        \texttt{elm\_Latn} & Eleme & L \\
        \texttt{emp\_Latn} & Northern Emberá & M \\
        \texttt{enb\_Latn} & Markweeta & M \\
        \texttt{eng\_Latn} & English & H \\
        \texttt{enx\_Latn} & Enxet & M \\
        \texttt{epo\_Latn} & Esperanto & H \\
        \texttt{ese\_Latn} & Ese Ejja & M \\
        \texttt{ess\_Latn} & Central Siberian Yupik & H \\
        \texttt{esu\_Latn} & Central Yupik & L \\
        \texttt{eto\_Latn} & Eton (Cameroon) & L \\
        \texttt{ets\_Latn} & Yekhee & L \\
        \texttt{etu\_Latn} & Ejagham & L \\
        \texttt{eus\_Latn} & Basque & H \\
        \texttt{evn\_Cyrl} & Evenki & L \\
        \texttt{ewe\_Latn} & Ewe & H \\
        \texttt{ewo\_Latn} & Ewondo & M \\
        \texttt{eyo\_Latn} & Keiyo & L \\
        \texttt{eza\_Latn} & Ezaa & M \\
        \texttt{fal\_Latn} & South Fali & M \\
        \texttt{fan\_Latn} & Fang (Equatorial Guinea) & M \\
        \texttt{fao\_Latn} & Faroese & H \\
        \texttt{far\_Latn} & Fataleka & H \\
        \texttt{fas\_Arab} & Persian & H \\
        \texttt{fat\_Latn} & Fanti & L \\
        \texttt{fia\_Latn} & Nobiin & L \\
        \texttt{fij\_Latn} & Fijian & M \\
        \texttt{fil\_Latn} & Filipino & H \\
        \texttt{fin\_Latn} & Finnish & H \\
        \texttt{fip\_Latn} & Fipa & L \\
        \texttt{fkk\_Latn} & Kirya-Konz\textschwa{}l & L \\
        \texttt{flr\_Latn} & Fuliiru & H \\
        \texttt{fmp\_Latn} & Fe'fe' & L \\
        \texttt{fmu\_Deva} & Far Western Muria & L \\
        \texttt{fon\_Latn} & Fon & M \\
        \texttt{fra\_Latn} & French & H \\
        \texttt{frd\_Latn} & Fordata & H \\
        \texttt{fry\_Latn} & Western Frisian & L \\
        \texttt{fub\_Latn} & Adamawa Fulfulde & M \\
        \texttt{fuc\_Latn} & Pulaar & L \\
        \texttt{fue\_Latn} & Borgu Fulfulde & L \\
        \texttt{ful\_Latn} & Fulah & H \\
        \texttt{fuq\_Latn} & Central-Eastern Niger Fulfulde & L \\
        \texttt{fuv\_Latn} & Nigerian Fulfulde & L \\
        \texttt{gag\_Cyrl} & Gagauz & M \\
        \texttt{gag\_Latn} & Gagauz & M \\
        \texttt{gai\_Latn} & Borei & H \\
        \texttt{gam\_Latn} & Kandawo & M \\
        \texttt{gau\_Telu} & Mudhili Gadaba & M \\
        \texttt{gbi\_Latn} & Galela & M \\
        \texttt{gbk\_Deva} & Gaddi & M \\
        \texttt{gbm\_Deva} & Garhwali & M \\
        \texttt{gbo\_Latn} & Northern Grebo & M \\
        \texttt{gbr\_Latn} & Gbagyi & L \\
        \texttt{gby\_Latn} & Gbari & L \\
        \texttt{gcc\_Latn} & Mali & L \\
        \texttt{gde\_Latn} & Gude & M \\
        \texttt{gdf\_Latn} & Guduf-Gava & L \\
        \texttt{geb\_Latn} & Kire & L \\
        \texttt{gej\_Latn} & Gen & M \\
        \texttt{ges\_Latn} & Geser-Gorom & L \\
        \texttt{ggg\_Arab} & Gurgula & L \\
        \texttt{gid\_Latn} & Gidar & L \\
        \texttt{gig\_Arab} & Goaria & L \\
        \texttt{gil\_Latn} & Gilbertese & M \\
        \texttt{giz\_Latn} & South Giziga & L \\
        \texttt{gjk\_Arab} & Kachi Koli & M \\
        \texttt{gjn\_Latn} & Gonja & M \\
        \texttt{gju\_Arab} & Gujari & L \\
        \texttt{gkn\_Latn} & Gokana & M \\
        \texttt{gld\_Cyrl} & Nanai & L \\
        \texttt{gle\_Latn} & Irish & H \\
        \texttt{glg\_Latn} & Galician & H \\
        \texttt{glk\_Arab} & Gilaki & L \\
        \texttt{glv\_Latn} & Manx & L \\
        \texttt{glw\_Latn} & Glavda & L \\
        \texttt{gmv\_Latn} & Gamo & M \\
        \texttt{gna\_Latn} & Kaansa & M \\
        \texttt{gnd\_Latn} & Zulgo-Gemzek & M \\
        \texttt{gng\_Latn} & Ngangam & M \\
        \texttt{gof\_Latn} & Gofa & M \\
        \texttt{gog\_Latn} & Gogo & M \\
        \texttt{gol\_Latn} & Gola & L \\
        \texttt{gom\_Deva} & Goan Konkani & L \\
        \texttt{gor\_Latn} & Gorontalo & M \\
        \texttt{gqr\_Latn} & Gor & M \\
        \texttt{grc\_Grek} & Ancient Greek (to 1453) & M \\
        \texttt{gri\_Latn} & Ghari & H \\
        \texttt{grn\_Latn} & Guarani & M \\
        \texttt{grt\_Beng} & Garo & M \\
        \texttt{gsl\_Latn} & Gusilay & L \\
        \texttt{gso\_Latn} & Southwest Gbaya & M \\
        \texttt{gub\_Latn} & Guajajára & H \\
        \texttt{guc\_Latn} & Wayuu & M \\
        \texttt{gud\_Latn} & Yocoboué Dida & M \\
        \texttt{gug\_Latn} & Paraguayan Guaraní & M \\
        \texttt{guh\_Latn} & Guahibo & H \\
        \texttt{gui\_Latn} & Eastern Bolivian Guaraní & L \\
        \texttt{guj\_Gujr} & Gujarati & H \\
        \texttt{guk\_Ethi} & Gumuz & M \\
        \texttt{gum\_Latn} & Guambiano & M \\
        \texttt{guo\_Latn} & Guayabero & M \\
        \texttt{guq\_Latn} & Aché & M \\
        \texttt{gur\_Latn} & Farefare & M \\
        \texttt{guu\_Latn} & Yanomamö & M \\
        \texttt{gux\_Latn} & Gourmanchéma & M \\
        \texttt{guz\_Latn} & Gusii & L \\
        \texttt{gvc\_Latn} & Guanano & M \\
        \texttt{gvl\_Latn} & Gulay & M \\
        \texttt{gwc\_Arab} & Gawri & L \\
        \texttt{gwe\_Latn} & Gweno & L \\
        \texttt{gwi\_Latn} & Gwich'in & M \\
        \texttt{gwr\_Latn} & Gwere & H \\
        \texttt{gwt\_Arab} & Gawar-Bati & L \\
        \texttt{gym\_Latn} & Ngäbere & M \\
        \texttt{gyr\_Latn} & Guarayu & M \\
        \bottomrule
    \end{tabular}
    \quad
    \begin{tabular}{llc}
        \toprule
        \textbf{Code} & \textbf{Name} & \textbf{Res}\\
        \midrule
        \texttt{gyz\_Latn} & Geji & L \\
        \texttt{had\_Latn} & Hatam & M \\
        \texttt{hag\_Latn} & Hanga & M \\
        \texttt{hah\_Latn} & Hahon & L \\
        \texttt{hak\_Latn} & Hakka Chinese & M \\
        \texttt{hao\_Latn} & Hakö & L \\
        \texttt{hap\_Latn} & Hupla & H \\
        \texttt{hat\_Latn} & Haitian & H \\
        \texttt{hau\_Latn} & Hausa & H \\
        \texttt{haw\_Latn} & Hawaiian & L \\
        \texttt{hay\_Latn} & Haya & M \\
        \texttt{hbb\_Latn} & Huba & L \\
        \texttt{hch\_Latn} & Huichol & L \\
        \texttt{heb\_Hebr} & Hebrew & H \\
        \texttt{heh\_Latn} & Hehe & M \\
        \texttt{her\_Latn} & Herero & L \\
        \texttt{hia\_Latn} & Lamang & L \\
        \texttt{hif\_Latn} & Fiji Hindi & M \\
        \texttt{hig\_Latn} & Kamwe & M \\
        \texttt{hil\_Latn} & Hiligaynon & M \\
        \texttt{hin\_Deva} & Hindi & H \\
        \texttt{hkk\_Latn} & Hunjara-Kaina Ke & L \\
        \texttt{hla\_Latn} & Halia & L \\
        \texttt{hlb\_Deva} & Halbi & M \\
        \texttt{hlt\_Latn} & Matu Chin & M \\
        \texttt{hne\_Deva} & Chhattisgarhi & H \\
        \texttt{hnn\_Latn} & Hanunoo & M \\
        \texttt{hno\_Arab} & Northern Hindko & M \\
        \texttt{hns\_Latn} & Caribbean Hindustani & M \\
        \texttt{hoc\_Orya} & Ho & H \\
        \texttt{hrv\_Latn} & Croatian & H \\
        \texttt{hsb\_Latn} & Upper Sorbian & L \\
        \texttt{hto\_Latn} & Minica Huitoto & M \\
        \texttt{hub\_Latn} & Huambisa & M \\
        \texttt{hue\_Latn} & San Francisco Del Mar Huave & L \\
        \texttt{hui\_Latn} & Huli & M \\
        \texttt{hul\_Latn} & Hula & L \\
        \texttt{hun\_Latn} & Hungarian & H \\
        \texttt{hus\_Latn} & Huastec & H \\
        \texttt{huu\_Latn} & Murui Huitoto & H \\
        \texttt{huv\_Latn} & San Mateo Del Mar Huave & M \\
        \texttt{hux\_Latn} & Nüpode Huitoto & L \\
        \texttt{hvn\_Latn} & Sabu & L \\
        \texttt{hwc\_Latn} & Hawai'i Creole English & M \\
        \texttt{hwo\_Latn} & Hwana & L \\
        \texttt{hye\_Armn} & Armenian & H \\
        \texttt{hyw\_Armn} & Western Armenian & M \\
        \texttt{iba\_Latn} & Iban & M \\
        \texttt{ibb\_Latn} & Ibibio & L \\
        \texttt{ibo\_Latn} & Igbo & H \\
        \texttt{icr\_Latn} & Islander Creole English & M \\
        \texttt{ida\_Latn} & Idakho-Isukha-Tiriki & L \\
        \texttt{idd\_Latn} & Ede Idaca & M \\
        \texttt{idu\_Latn} & Idoma & L \\
        \texttt{ifa\_Latn} & Amganad Ifugao & M \\
        \texttt{ifb\_Latn} & Batad Ifugao & M \\
        \texttt{ife\_Latn} & Ifè & H \\
        \texttt{ifk\_Latn} & Tuwali Ifugao & M \\
        \texttt{ifu\_Latn} & Mayoyao Ifugao & M \\
        \texttt{ify\_Latn} & Keley-I Kallahan & M \\
        \texttt{igl\_Latn} & Igala & L \\
        \texttt{ign\_Latn} & Ignaciano & M \\
        \texttt{ijc\_Latn} & Izon & L \\
        \texttt{ijn\_Latn} & Kalabari & L \\
        \texttt{ikk\_Latn} & Ika & M \\
        \texttt{ikw\_Latn} & Ikwere & L \\
        \texttt{ilb\_Latn} & Ila & M \\
        \texttt{ilo\_Latn} & Iloko & M \\
        \texttt{imo\_Latn} & Imbongu & L \\
        \texttt{ina\_Latn} & Interlingua & L \\
        \texttt{inb\_Latn} & Inga & M \\
        \texttt{ind\_Latn} & Indonesian & H \\
        \texttt{iou\_Latn} & Tuma-Irumu & M \\
        \texttt{ipi\_Latn} & Ipili & M \\
        \texttt{ipk\_Latn} & Inupiaq & L \\
        \texttt{iqw\_Latn} & Ikwo & M \\
        \texttt{iri\_Latn} & Rigwe & M \\
        \texttt{irk\_Latn} & Iraqw & M \\
        \texttt{ish\_Latn} & Esan & L \\
        \texttt{isl\_Latn} & Icelandic & H \\
        \texttt{iso\_Latn} & Isoko & L \\
        \texttt{ita\_Latn} & Italian & H \\
        \texttt{itl\_Cyrl} & Itelmen & L \\
        \texttt{its\_Latn} & Isekiri & L \\
        \texttt{itv\_Latn} & Itawit & M \\
        \texttt{itw\_Latn} & Ito & L \\
        \texttt{itz\_Latn} & Itzá & L \\
        \texttt{ixl\_Latn} & Ixil & H \\
        \texttt{izr\_Latn} & Izere & M \\
        \texttt{izz\_Latn} & Izii & M \\
        \texttt{jac\_Latn} & Popti' & M \\
        \texttt{jal\_Latn} & Yalahatan & L \\
        \texttt{jam\_Latn} & Jamaican Creole English & M \\
        \texttt{jav\_Latn} & Javanese & H \\
        \texttt{jax\_Latn} & Jambi Malay & L \\
        \texttt{jbu\_Latn} & Jukun Takum & M \\
        \texttt{jen\_Latn} & Dza & L \\
        \texttt{jic\_Latn} & Tol & M \\
        \texttt{jiv\_Latn} & Shuar & M \\
        \texttt{jmc\_Latn} & Machame & M \\
        \texttt{jmd\_Latn} & Yamdena & L \\
        \texttt{jmx\_Latn} & Western Juxtlahuaca Mixtec & L \\
        \texttt{jpn\_Jpan} & Japanese & H \\
        \texttt{jqr\_Latn} & Jaqaru & L \\
        \texttt{juk\_Latn} & Wapan & L \\
        \texttt{jun\_Orya} & Juang & H \\
        \texttt{juo\_Latn} & Jiba & L \\
        \texttt{jvn\_Latn} & Caribbean Javanese & M \\
        \texttt{kaa\_Cyrl} & Kara-Kalpak & M \\
        \texttt{kab\_Latn} & Kabyle & H \\
        \texttt{kac\_Latn} & Kachin & M \\
        \texttt{kai\_Latn} & Karekare & L \\
        \texttt{kaj\_Latn} & Jju & L \\
        \texttt{kak\_Latn} & Kalanguya & M \\
        \texttt{kam\_Latn} & Kamba (Kenya) & M \\
        \texttt{kan\_Knda} & Kannada & H \\
        \texttt{kao\_Latn} & Xaasongaxango & M \\
        \texttt{kaq\_Latn} & Capanahua & H \\
        \texttt{kas\_Arab} & Kashmiri & L \\
        \texttt{kat\_Geor} & Georgian & H \\
        \texttt{kay\_Latn} & Kamayurá & L \\
        \texttt{kaz\_Cyrl} & Kazakh & H \\
        \texttt{kbd\_Cyrl} & Kabardian & H \\
        \texttt{kbl\_Latn} & Kanembu & L \\
        \texttt{kbo\_Latn} & Keliko & H \\
        \texttt{kbp\_Latn} & Kabiyè & H \\
        \texttt{kbq\_Latn} & Kamano & M \\
        \texttt{kbr\_Latn} & Kafa & M \\
        \texttt{kbt\_Latn} & Abadi & L \\
        \texttt{kby\_Latn} & Manga Kanuri & L \\
        \texttt{kca\_Cyrl} & Khanty & L \\
        \texttt{kcg\_Latn} & Tyap & M \\
        \texttt{kcn\_Latn} & Nubi & L \\
        \texttt{kcq\_Latn} & Kamo & L \\
        \texttt{kdc\_Latn} & Kutu & M \\
        \texttt{kde\_Latn} & Makonde & M \\
        \texttt{kdh\_Latn} & Tem & M \\
        \texttt{kdi\_Latn} & Kumam & M \\
        \texttt{kdj\_Latn} & Karamojong & M \\
        \texttt{kdl\_Latn} & Tsikimba & M \\
        \texttt{kdn\_Latn} & Kunda & M \\
        \texttt{kdt\_Khmr} & Kuy & M \\
        \texttt{kea\_Latn} & Kabuverdianu & M \\
        \texttt{kek\_Latn} & Kekchí & M \\
        \texttt{ken\_Latn} & Kenyang & M \\
        \texttt{keo\_Latn} & Kakwa & M \\
        \texttt{ker\_Latn} & Kera & M \\
        \texttt{keu\_Latn} & Akebu & L \\
        \texttt{key\_Telu} & Kupia & M \\
        \texttt{kez\_Latn} & Kukele & M \\
        \bottomrule
    \end{tabular}
    \quad
    \begin{tabular}{llc}
        \toprule
        \textbf{Code} & \textbf{Name} & \textbf{Res}\\
        \midrule
        \texttt{kfb\_Deva} & Northwestern Kolami & M \\
        \texttt{kff\_Telu} & Koya & M \\
        \texttt{kfk\_Deva} & Kinnauri & L \\
        \texttt{kfq\_Deva} & Korku & L \\
        \texttt{kfr\_Gujr} & Kachhi & L \\
        \texttt{kfw\_Latn} & Kharam Naga & M \\
        \texttt{kfx\_Deva} & Kullu Pahari & M \\
        \texttt{kha\_Latn} & Khasi & L \\
        \texttt{khg\_Tibt} & Khams Tibetan & M \\
        \texttt{khk\_Cyrl} & Halh Mongolian & M \\
        \texttt{khm\_Khmr} & Khmer & H \\
        \texttt{khq\_Latn} & Koyra Chiini Songhay & M \\
        \texttt{khw\_Arab} & Khowar & M \\
        \texttt{kia\_Latn} & Kim & M \\
        \texttt{kij\_Latn} & Kilivila & M \\
        \texttt{kik\_Latn} & Kikuyu & M \\
        \texttt{kin\_Latn} & Kinyarwanda & H \\
        \texttt{kir\_Cyrl} & Kirghiz & H \\
        \texttt{kix\_Latn} & Khiamniungan Naga & L \\
        \texttt{kjb\_Latn} & Q'anjob'al & M \\
        \texttt{kjc\_Latn} & Coastal Konjo & L \\
        \texttt{kje\_Latn} & Kisar & H \\
        \texttt{kjg\_Latn} & Khmu & L \\
        \texttt{kjh\_Cyrl} & Khakas & M \\
        \texttt{kjk\_Latn} & Highland Konjo & L \\
        \texttt{kki\_Latn} & Kagulu & M \\
        \texttt{kkj\_Latn} & Kako & M \\
        \texttt{kle\_Deva} & Kulung (Nepal) & M \\
        \texttt{kln\_Latn} & Kalenjin & M \\
        \texttt{kls\_Latn} & Kalasha & L \\
        \texttt{klu\_Latn} & Klao & M \\
        \texttt{klv\_Latn} & Maskelynes & M \\
        \texttt{klw\_Latn} & Tado & L \\
        \texttt{kma\_Latn} & Konni & M \\
        \texttt{kmd\_Latn} & Majukayang Kalinga & M \\
        \texttt{kml\_Latn} & Tanudan Kalinga & M \\
        \texttt{kmr\_Arab} & Northern Kurdish & M \\
        \texttt{kmr\_Cyrl} & Northern Kurdish & M \\
        \texttt{kmr\_Latn} & Northern Kurdish & H \\
        \texttt{kmu\_Latn} & Kanite & H \\
        \texttt{kmy\_Latn} & Koma & L \\
        \texttt{kna\_Latn} & Dera (Nigeria) & L \\
        \texttt{knb\_Latn} & Lubuagan Kalinga & M \\
        \texttt{knc\_Latn} & Central Kanuri & L \\
        \texttt{kne\_Latn} & Kankanaey & M \\
        \texttt{knf\_Latn} & Mankanya & M \\
        \texttt{knj\_Latn} & Western Kanjobal & M \\
        \texttt{knk\_Latn} & Kuranko & M \\
        \texttt{knn\_Deva} & Konkani (individual language) & L \\
        \texttt{kno\_Latn} & Kono (Sierra Leone) & M \\
        \texttt{kog\_Latn} & Cogui & H \\
        \texttt{kol\_Latn} & Kol (Papua New Guinea) & L \\
        \texttt{koo\_Latn} & Konzo & M \\
        \texttt{kor\_Hang} & Korean & H \\
        \texttt{kpo\_Latn} & Ikposo & L \\
        \texttt{kpq\_Latn} & Korupun-Sela & H \\
        \texttt{kps\_Latn} & Tehit & L \\
        \texttt{kpv\_Cyrl} & Komi-Zyrian & M \\
        \texttt{kpy\_Cyrl} & Koryak & L \\
        \texttt{kpz\_Latn} & Kupsabiny & M \\
        \texttt{kqe\_Latn} & Kalagan & H \\
        \texttt{kqo\_Latn} & Eastern Krahn & L \\
        \texttt{kqp\_Latn} & Kimré & M \\
        \texttt{kqr\_Latn} & Kimaragang & H \\
        \texttt{kqy\_Ethi} & Koorete & M \\
        \texttt{krc\_Cyrl} & Karachay-Balkar & M \\
        \texttt{kri\_Latn} & Krio & M \\
        \texttt{krj\_Latn} & Kinaray-A & M \\
        \texttt{krl\_Latn} & Karelian & M \\
        \texttt{krr\_Khmr} & Krung & L \\
        \texttt{krs\_Latn} & Gbaya (Sudan) & H \\
        \texttt{kru\_Deva} & Kurukh & M \\
        \texttt{krx\_Latn} & Karon & L \\
        \texttt{ksb\_Latn} & Shambala & M \\
        \texttt{ksd\_Latn} & Kuanua & L \\
        \texttt{ksf\_Latn} & Bafia & M \\
        \texttt{ksr\_Latn} & Borong & H \\
        \texttt{kss\_Latn} & Southern Kisi & M \\
        \texttt{ksz\_Deva} & Kodaku & L \\
        \texttt{ktb\_Ethi} & Kambaata & M \\
        \texttt{ktj\_Latn} & Plapo Krumen & H \\
        \texttt{kto\_Latn} & Kuot & L \\
        \texttt{kua\_Latn} & Kuanyama & L \\
        \texttt{kub\_Latn} & Kutep & M \\
        \texttt{kue\_Latn} & Kuman (Papua New Guinea) & M \\
        \texttt{kuh\_Latn} & Kushi & L \\
        \texttt{kum\_Cyrl} & Kumyk & M \\
        \texttt{kur\_Arab} & Kurdish & M \\
        \texttt{kus\_Latn} & Kusaal & M \\
        \texttt{kvn\_Latn} & Border Kuna & H \\
        \texttt{kvw\_Latn} & Wersing & L \\
        \texttt{kvx\_Arab} & Parkari Koli & L \\
        \texttt{kwd\_Latn} & Kwaio & H \\
        \texttt{kwf\_Latn} & Kwara'ae & H \\
        \texttt{kwi\_Latn} & Awa-Cuaiquer & M \\
        \texttt{kwm\_Latn} & Kwambi & L \\
        \texttt{kxc\_Ethi} & Konso & M \\
        \texttt{kxf\_Latn} & Manumanaw Karen & M \\
        \texttt{kxm\_Thai} & Northern Khmer & M \\
        \texttt{kxp\_Arab} & Wadiyara Koli & M \\
        \texttt{kyb\_Latn} & Butbut Kalinga & M \\
        \texttt{kyc\_Latn} & Kyaka & M \\
        \texttt{kyf\_Latn} & Kouya & M \\
        \texttt{kyg\_Latn} & Keyagana & L \\
        \texttt{kyo\_Latn} & Kelon & L \\
        \texttt{kyq\_Latn} & Kenga & M \\
        \texttt{kyu\_Kali} & Western Kayah & M \\
        \texttt{kyx\_Latn} & Rapoisi & L \\
        \texttt{kyz\_Latn} & Kayabí & H \\
        \texttt{kzf\_Latn} & Da'a Kaili & H \\
        \texttt{kzi\_Latn} & Kelabit & L \\
        \texttt{lac\_Latn} & Lacandon & M \\
        \texttt{lag\_Latn} & Rangi & L \\
        \texttt{laj\_Latn} & Lango (Uganda) & M \\
        \texttt{lam\_Latn} & Lamba & M \\
        \texttt{lao\_Laoo} & Lao & H \\
        \texttt{las\_Latn} & Lama (Togo) & M \\
        \texttt{lat\_Latn} & Latin & M \\
        \texttt{lav\_Latn} & Latvian & H \\
        \texttt{law\_Latn} & Lauje & H \\
        \texttt{lbj\_Tibt} & Ladakhi & L \\
        \texttt{lbw\_Latn} & Tolaki & M \\
        \texttt{lcm\_Latn} & Tungag & L \\
        \texttt{lcp\_Thai} & Western Lawa & M \\
        \texttt{ldb\_Latn} & D\~uya & L \\
        \texttt{led\_Latn} & Lendu & L \\
        \texttt{lee\_Latn} & Lyélé & M \\
        \texttt{lef\_Latn} & Lelemi & M \\
        \texttt{lem\_Latn} & Nomaande & M \\
        \texttt{lew\_Latn} & Ledo Kaili & H \\
        \texttt{lex\_Latn} & Luang & H \\
        \texttt{lgg\_Latn} & Lugbara & M \\
        \texttt{lgl\_Latn} & Wala & M \\
        \texttt{lhu\_Latn} & Lahu & M \\
        \texttt{lia\_Latn} & West-Central Limba & M \\
        \texttt{lid\_Latn} & Nyindrou & H \\
        \texttt{lif\_Deva} & Limbu & M \\
        \texttt{lij\_Latn} & Ligurian & M \\
        \texttt{lin\_Latn} & Lingala & H \\
        \texttt{lip\_Latn} & Sekpele & M \\
        \texttt{lir\_Latn} & Liberian English & L \\
        \texttt{lis\_Lisu} & Lisu & M \\
        \texttt{lit\_Latn} & Lithuanian & H \\
        \texttt{lje\_Latn} & Rampi & M \\
        \texttt{ljp\_Latn} & Lampung Api & M \\
        \texttt{lkb\_Latn} & Kabras & L \\
        \texttt{lke\_Latn} & Kenyi & L \\
        \texttt{lla\_Latn} & Lala-Roba & L \\
        \texttt{lld\_Latn\_gherd} & Ladin (Gherdëina) & L \\
        \texttt{lld\_Latn\_valbadia} & Ladin (Val Badia) & L \\
        \bottomrule
    \end{tabular}
    \end{adjustbox}
}

{
    \tiny\renewcommand{\arraystretch}{.8}
    \begin{adjustbox}{max width=\textwidth}
    \begin{tabular}{llc}
        \toprule
        \textbf{Code} & \textbf{Name} & \textbf{Res}\\
        \midrule
        \texttt{llg\_Latn} & Lole & L \\
        \texttt{lln\_Latn} & Lele (Chad) & H \\
        \texttt{lme\_Latn} & Pévé & M \\
        \texttt{lnd\_Latn} & Lundayeh & M \\
        \texttt{lns\_Latn} & Lamnso' & M \\
        \texttt{lnu\_Latn} & Longuda & L \\
        \texttt{loa\_Latn} & Loloda & L \\
        \texttt{lob\_Latn} & Lobi & M \\
        \texttt{lok\_Latn} & Loko & M \\
        \texttt{lom\_Latn} & Loma (Liberia) & M \\
        \texttt{lon\_Latn} & Malawi Lomwe & M \\
        \texttt{loq\_Latn} & Lobala & M \\
        \texttt{lrk\_Arab} & Loarki & L \\
        \texttt{lsi\_Latn} & Lashi & M \\
        \texttt{lsm\_Latn} & Saamia & M \\
        \texttt{lss\_Arab} & Lasi & L \\
        \texttt{ltg\_Latn} & Latgalian & L \\
        \texttt{lth\_Latn} & Thur & L \\
        \texttt{lto\_Latn} & Tsotso & L \\
        \texttt{ltz\_Latn} & Luxembourgish & L \\
        \texttt{lua\_Latn} & Luba-Lulua & L \\
        \texttt{luc\_Latn} & Aringa & M \\
        \texttt{lug\_Latn} & Ganda & H \\
        \texttt{luo\_Latn} & Luo (Kenya and Tanzania) & H \\
        \texttt{lus\_Latn} & Lushai & L \\
        \texttt{lwg\_Latn} & Wanga & L \\
        \texttt{lwo\_Latn} & Luwo & M \\
        \texttt{lww\_Latn} & Lewo & M \\
        \texttt{lzz\_Latn} & Laz & L \\
        \texttt{maa\_Latn} & San Jerónimo Tecóatl Mazatec & M \\
        \texttt{mab\_Latn} & Yutanduchi Mixtec & L \\
        \texttt{mad\_Latn} & Madurese & M \\
        \texttt{maf\_Latn} & Mafa & L \\
        \texttt{mag\_Deva} & Magahi & M \\
        \texttt{mah\_Latn} & Marshallese & M \\
        \texttt{mai\_Deva} & Maithili & M \\
        \texttt{maj\_Latn} & Jalapa De Díaz Mazatec & M \\
        \texttt{mak\_Latn} & Makasar & M \\
        \texttt{mal\_Mlym} & Malayalam & H \\
        \texttt{mam\_Latn} & Mam & H \\
        \texttt{maq\_Latn} & Chiquihuitlán Mazatec & L \\
        \texttt{mar\_Deva} & Marathi & H \\
        \texttt{mau\_Latn} & Huautla Mazatec & L \\
        \texttt{maw\_Latn} & Mampruli & M \\
        \texttt{max\_Latn} & North Moluccan Malay & L \\
        \texttt{maz\_Latn} & Central Mazahua & M \\
        \texttt{mbb\_Latn} & Western Bukidnon Manobo & M \\
        \texttt{mbc\_Latn} & Macushi & H \\
        \texttt{mbh\_Latn} & Mangseng & M \\
        \texttt{mbj\_Latn} & Nadëb & M \\
        \texttt{mbt\_Latn} & Matigsalug Manobo & M \\
        \texttt{mbu\_Latn} & Mbula-Bwazza & M \\
        \texttt{mca\_Latn} & Maca & M \\
        \texttt{mcb\_Latn} & Machiguenga & M \\
        \texttt{mcd\_Latn} & Sharanahua & H \\
        \texttt{mcf\_Latn} & Matsés & L \\
        \texttt{mco\_Latn} & Coatlán Mixe & M \\
        \texttt{mcp\_Latn} & Makaa & M \\
        \texttt{mcq\_Latn} & Ese & M \\
        \texttt{mcu\_Latn} & Cameroon Mambila & M \\
        \texttt{mcx\_Latn} & Mpiemo & L \\
        \texttt{mda\_Latn} & Mada (Nigeria) & M \\
        \texttt{mdd\_Latn} & Mbum & L \\
        \texttt{mdv\_Latn} & Santa Lucía Monteverde Mixtec & L \\
        \texttt{mdy\_Ethi} & Male (Ethiopia) & M \\
        \texttt{med\_Latn} & Melpa & M \\
        \texttt{mee\_Latn} & Mengen & L \\
        \texttt{meh\_Latn} & Southwestern Tlaxiaco Mixtec & L \\
        \texttt{mej\_Latn} & Meyah & H \\
        \texttt{mek\_Latn} & Mekeo & L \\
        \texttt{mel\_Latn} & Central Melanau & L \\
        \texttt{men\_Latn} & Mende (Sierra Leone) & M \\
        \texttt{meq\_Latn} & Merey & M \\
        \texttt{mer\_Latn} & Meru & L \\
        \texttt{met\_Latn} & Mato & M \\
        \texttt{meu\_Latn} & Motu & L \\
        \texttt{mev\_Latn} & Mano & M \\
        \texttt{mfe\_Latn} & Morisyen & M \\
        \texttt{mfh\_Latn} & Matal & M \\
        \texttt{mfi\_Latn} & Wandala & M \\
        \texttt{mfk\_Latn} & North Mofu & M \\
        \texttt{mfm\_Latn} & Marghi South & L \\
        \texttt{mfn\_Latn} & Cross River Mbembe & L \\
        \texttt{mfo\_Latn} & Mbe & L \\
        \texttt{mfq\_Latn} & Moba & M \\
        \texttt{mfv\_Latn} & Mandjak & L \\
        \texttt{mfy\_Latn} & Mayo & M \\
        \texttt{mfz\_Latn} & Mabaan & M \\
        \texttt{mgd\_Latn} & Moru & M \\
        \texttt{mge\_Latn} & Mango & M \\
        \texttt{mgg\_Latn} & Mpumpong & L \\
        \texttt{mgh\_Latn} & Makhuwa-Meetto & M \\
        \texttt{mgi\_Latn} & Lijili & L \\
        \texttt{mgo\_Latn} & Meta' & M \\
        \texttt{mhi\_Latn} & Ma'di & M \\
        \texttt{mhk\_Latn} & Mungaka & L \\
        \texttt{mhr\_Cyrl} & Eastern Mari & H \\
        \texttt{mhu\_Latn} & Digaro-Mishmi & L \\
        \texttt{mhx\_Latn} & Maru & H \\
        \texttt{mhy\_Latn} & Ma'anyan & M \\
        \texttt{mib\_Latn} & Atatláhuca Mixtec & H \\
        \texttt{mie\_Latn} & Ocotepec Mixtec & M \\
        \texttt{mif\_Latn} & Mofu-Gudur & M \\
        \texttt{mig\_Latn} & San Miguel El Grande Mixtec & L \\
        \texttt{mih\_Latn} & Chayuco Mixtec & H \\
        \texttt{mil\_Latn} & Peñoles Mixtec & H \\
        \texttt{mim\_Latn} & Alacatlatzala Mixtec & H \\
        \texttt{min\_Latn} & Minangkabau & M \\
        \texttt{mio\_Latn} & Pinotepa Nacional Mixtec & M \\
        \texttt{mip\_Latn} & Apasco-Apoala Mixtec & M \\
        \texttt{miq\_Latn} & Mískito & M \\
        \texttt{mit\_Latn} & Southern Puebla Mixtec & M \\
        \texttt{miu\_Latn} & Cacaloxtepec Mixtec & L \\
        \texttt{miy\_Latn} & Ayutla Mixtec & M \\
        \texttt{miz\_Latn} & Coatzospan Mixtec & M \\
        \texttt{mjl\_Deva} & Mandeali & M \\
        \texttt{mjv\_Mlym} & Mannan & M \\
        \texttt{mkd\_Cyrl} & Macedonian & H \\
        \texttt{mkf\_Latn} & Miya & L \\
        \texttt{mki\_Arab} & Dhatki & L \\
        \texttt{mkl\_Latn} & Mokole & M \\
        \texttt{mkn\_Latn} & Kupang Malay & L \\
        \texttt{mlg\_Latn} & Malagasy & H \\
        \texttt{mlq\_Latn} & Western Maninkakan & L \\
        \texttt{mlt\_Latn} & Maltese & H \\
        \texttt{mmc\_Latn} & Michoacán Mazahua & L \\
        \texttt{mmg\_Latn} & North Ambrym & L \\
        \texttt{mnb\_Latn} & Muna & M \\
        \texttt{mne\_Latn} & Naba & L \\
        \texttt{mnf\_Latn} & Mundani & M \\
        \texttt{mni\_Beng} & Manipuri & L \\
        \texttt{mnk\_Latn} & Mandinka & M \\
        \texttt{mnw\_Mymr} & Mon & M \\
        \texttt{mnx\_Latn} & Manikion & H \\
        \texttt{moa\_Latn} & Mwan & H \\
        \texttt{mog\_Latn} & Mongondow & M \\
        \texttt{mon\_Cyrl} & Mongolian & M \\
        \texttt{mop\_Latn} & Mopán Maya & M \\
        \texttt{mor\_Latn} & Moro & M \\
        \texttt{mos\_Latn} & Mossi & M \\
        \texttt{mox\_Latn} & Molima & M \\
        \texttt{moz\_Latn} & Mukulu & M \\
        \texttt{mpg\_Latn} & Marba & M \\
        \texttt{mpm\_Latn} & Yosondúa Mixtec & H \\
        \texttt{mpp\_Latn} & Migabac & M \\
        \texttt{mpx\_Latn} & Misima-Panaeati & M \\
        \texttt{mqb\_Latn} & Mbuko & M \\
        \texttt{mqf\_Latn} & Momuna & H \\
        \texttt{mqj\_Latn} & Mamasa & M \\
        \texttt{mqn\_Latn} & Moronene & M \\
        \bottomrule
    \end{tabular}
    \quad
    \begin{tabular}{llc}
        \toprule
        \textbf{Code} & \textbf{Name} & \textbf{Res}\\
        \midrule
        \texttt{mqy\_Latn} & Manggarai & L \\
        \texttt{mri\_Latn} & Maori & M \\
        \texttt{mrj\_Cyrl} & Western Mari & M \\
        \texttt{mrr\_Deva} & Maria (India) & L \\
        \texttt{mrt\_Latn} & Marghi Central & L \\
        \texttt{mrw\_Latn} & Maranao & M \\
        \texttt{msh\_Latn} & Masikoro Malagasy & L \\
        \texttt{msi\_Latn} & Sabah Malay & M \\
        \texttt{msw\_Latn} & Mansoanka & L \\
        \texttt{msy\_Latn} & Aruamu & M \\
        \texttt{mtd\_Latn} & Mualang & L \\
        \texttt{mtj\_Latn} & Moskona & H \\
        \texttt{mto\_Latn} & Totontepec Mixe & H \\
        \texttt{mtr\_Deva} & Mewari & L \\
        \texttt{mtu\_Latn} & Tututepec Mixtec & L \\
        \texttt{mtx\_Latn} & Tidaá Mixtec & L \\
        \texttt{mua\_Latn} & Mundang & L \\
        \texttt{mug\_Latn} & Musgu & L \\
        \texttt{muh\_Latn} & Mündü & M \\
        \texttt{mui\_Latn} & Musi & L \\
        \texttt{mup\_Deva} & Malvi & M \\
        \texttt{mur\_Latn} & Murle & M \\
        \texttt{muv\_Mlym} & Muthuvan & M \\
        \texttt{muy\_Latn} & Muyang & M \\
        \texttt{mve\_Arab} & Marwari (Pakistan) & L \\
        \texttt{mvp\_Latn} & Duri & M \\
        \texttt{mvy\_Arab} & Indus Kohistani & M \\
        \texttt{mwq\_Latn} & Mün Chin & M \\
        \texttt{mwv\_Latn} & Mentawai & M \\
        \texttt{mxb\_Latn} & Tezoatlán Mixtec & H \\
        \texttt{mxq\_Latn} & Juquila Mixe & M \\
        \texttt{mxs\_Latn} & Huitepec Mixtec & L \\
        \texttt{mxt\_Latn} & Jamiltepec Mixtec & M \\
        \texttt{mxu\_Latn} & Mada (Cameroon) & L \\
        \texttt{mxv\_Latn} & Metlatónoc Mixtec & M \\
        \texttt{mxy\_Latn} & Southeastern Nochixtlán Mixtec & L \\
        \texttt{mya\_Mymr} & Burmese & H \\
        \texttt{myb\_Latn} & Mbay & M \\
        \texttt{myk\_Latn} & Mamara Senoufo & M \\
        \texttt{myv\_Cyrl} & Erzya & M \\
        \texttt{myx\_Latn} & Masaaba & M \\
        \texttt{myy\_Latn} & Macuna & H \\
        \texttt{mza\_Latn} & Santa María Zacatepec Mixtec & M \\
        \texttt{mzi\_Latn} & Ixcatlán Mazatec & M \\
        \texttt{mzj\_Latn} & Manya & M \\
        \texttt{mzk\_Latn} & Nigeria Mambila & M \\
        \texttt{mzl\_Latn} & Mazatlán Mixe & L \\
        \texttt{mzm\_Latn} & Mumuye & M \\
        \texttt{mzw\_Latn} & Deg & M \\
        \texttt{nab\_Latn} & Southern Nambikuára & M \\
        \texttt{nag\_Latn} & Naga Pidgin & M \\
        \texttt{nal\_Latn} & Nalik & L \\
        \texttt{nan\_Latn} & Min Nan Chinese & M \\
        \texttt{nap\_Latn} & Neapolitan & L \\
        \texttt{nas\_Latn} & Naasioi & M \\
        \texttt{naw\_Latn} & Nawuri & M \\
        \texttt{nbh\_Latn} & Ngamo & L \\
        \texttt{nca\_Latn} & Iyo & M \\
        \texttt{ncf\_Latn} & Notsi & L \\
        \texttt{nch\_Latn} & Central Huasteca Nahuatl & M \\
        \texttt{ncj\_Latn} & Northern Puebla Nahuatl & M \\
        \texttt{ncl\_Latn} & Michoacán Nahuatl & M \\
        \texttt{nco\_Latn} & Sibe & L \\
        \texttt{ncu\_Latn} & Chumburung & M \\
        \texttt{ncx\_Latn} & Central Puebla Nahuatl & L \\
        \texttt{ndi\_Latn} & Samba Leko & L \\
        \texttt{ndj\_Latn} & Ndamba & M \\
        \texttt{ndo\_Latn} & Ndonga & L \\
        \texttt{ndp\_Latn} & Ndo & M \\
        \texttt{ndv\_Latn} & Ndut & M \\
        \texttt{ndy\_Latn} & Lutos & M \\
        \texttt{ndz\_Latn} & Ndogo & M \\
        \texttt{neb\_Latn} & Toura (Côte d'Ivoire) & M \\
        \texttt{nep\_Deva} & Nepali (macrolanguage) & M \\
        \texttt{new\_Deva} & Newari & M \\
        \texttt{nfa\_Latn} & Dhao & L \\
        \texttt{nfr\_Latn} & Nafaanra & M \\
        \texttt{nga\_Latn} & Ngbaka & M \\
        \texttt{ngi\_Latn} & Ngizim & L \\
        \texttt{ngl\_Latn} & Lomwe & M \\
        \texttt{ngp\_Latn} & Ngulu & M \\
        \texttt{ngu\_Latn} & Guerrero Nahuatl & M \\
        \texttt{nhe\_Latn} & Eastern Huasteca Nahuatl & M \\
        \texttt{nhg\_Latn} & Tetelcingo Nahuatl & L \\
        \texttt{nhi\_Latn} & Zacatlán-Ahuacatlán-Tepetzintla Nahuatl & M \\
        \texttt{nhn\_Latn} & Central Nahuatl & L \\
        \texttt{nhq\_Latn} & Huaxcaleca Nahuatl & L \\
        \texttt{nhu\_Latn} & Noone & M \\
        \texttt{nhw\_Latn} & Western Huasteca Nahuatl & M \\
        \texttt{nhx\_Latn} & Isthmus-Mecayapan Nahuatl & H \\
        \texttt{nhy\_Latn} & Northern Oaxaca Nahuatl & M \\
        \texttt{nia\_Latn} & Nias & M \\
        \texttt{nij\_Latn} & Ngaju & M \\
        \texttt{nim\_Latn} & Nilamba & M \\
        \texttt{nin\_Latn} & Ninzo & M \\
        \texttt{nja\_Latn} & Nzanyi & L \\
        \texttt{nko\_Latn} & Nkonya & M \\
        \texttt{nla\_Latn} & Ngombale & L \\
        \texttt{nlc\_Latn} & Nalca & H \\
        \texttt{nld\_Latn} & Dutch & H \\
        \texttt{nlg\_Latn} & Gela & H \\
        \texttt{nlk\_Latn} & Ninia Yali & H \\
        \texttt{nlv\_Latn} & Orizaba Nahuatl & L \\
        \texttt{nmg\_Latn} & Kwasio & L \\
        \texttt{nmz\_Latn} & Nawdm & M \\
        \texttt{nnb\_Latn} & Nande & M \\
        \texttt{nnh\_Latn} & Ngiemboon & M \\
        \texttt{nnq\_Latn} & Ngindo & M \\
        \texttt{nnw\_Latn} & Southern Nuni & M \\
        \texttt{noa\_Latn} & Woun Meu & M \\
        \texttt{nob\_Latn} & Norwegian Bokmål & H \\
        \texttt{nod\_Thai} & Northern Thai & M \\
        \texttt{noe\_Deva} & Nimadi & L \\
        \texttt{nog\_Cyrl} & Nogai & M \\
        \texttt{not\_Latn} & Nomatsiguenga & H \\
        \texttt{npl\_Latn} & Southeastern Puebla Nahuatl & M \\
        \texttt{npy\_Latn} & Napu & M \\
        \texttt{nso\_Latn} & Pedi & M \\
        \texttt{nst\_Latn} & Tase Naga & M \\
        \texttt{nsu\_Latn} & Sierra Negra Nahuatl & M \\
        \texttt{ntm\_Latn} & Nateni & M \\
        \texttt{ntr\_Latn} & Delo & M \\
        \texttt{nuj\_Latn} & Nyole & M \\
        \texttt{nup\_Latn} & Nupe-Nupe-Tako & L \\
        \texttt{nus\_Latn} & Nuer & M \\
        \texttt{nuz\_Latn} & Tlamacazapa Nahuatl & L \\
        \texttt{nwb\_Latn} & Nyabwa & M \\
        \texttt{nxq\_Latn} & Naxi & H \\
        \texttt{nya\_Latn} & Nyanja & M \\
        \texttt{nyf\_Latn} & Giryama & M \\
        \texttt{nyn\_Latn} & Nyankole & M \\
        \texttt{nyo\_Latn} & Nyoro & M \\
        \texttt{nyu\_Latn} & Nyungwe & L \\
        \texttt{nyy\_Latn} & Nyakyusa-Ngonde & M \\
        \texttt{nzi\_Latn} & Nzima & M \\
        \texttt{obo\_Latn} & Obo Manobo & H \\
        \texttt{oci\_Latn} & Occitan & M \\
        \texttt{odk\_Arab} & Od & M \\
        \texttt{odu\_Latn} & Odual & L \\
        \texttt{ogo\_Latn} & Khana & L \\
        \texttt{ojb\_Cans} & Northwestern Ojibwa & M \\
        \texttt{ojb\_Latn} & Northwestern Ojibwa & M \\
        \texttt{oku\_Latn} & Oku & M \\
        \texttt{old\_Latn} & Mochi & M \\
        \texttt{omw\_Latn} & South Tairora & H \\
        \texttt{onb\_Latn} & Lingao & M \\
        \texttt{ood\_Latn} & Tohono O'odham & M \\
        \texttt{orc\_Latn} & Orma & M \\
        \texttt{orm\_Latn} & Oromo & M \\
        \texttt{oru\_Arab} & Ormuri & M \\
        \bottomrule
    \end{tabular}
    \quad
    \begin{tabular}{llc}
        \toprule
        \textbf{Code} & \textbf{Name} & \textbf{Res}\\
        \midrule
        \texttt{ory\_Orya} & Odia & H \\
        \texttt{oss\_Cyrl} & Ossetian & M \\
        \texttt{ote\_Latn} & Mezquital Otomi & M \\
        \texttt{otq\_Latn} & Querétaro Otomi & M \\
        \texttt{ozm\_Latn} & Koonzime & M \\
        \texttt{pab\_Latn} & Parecís & M \\
        \texttt{pad\_Latn} & Paumarí & M \\
        \texttt{pag\_Latn} & Pangasinan & M \\
        \texttt{pam\_Latn} & Pampanga & M \\
        \texttt{pan\_Guru} & Panjabi & H \\
        \texttt{pao\_Latn} & Northern Paiute & M \\
        \texttt{pap\_Latn} & Papiamento & M \\
        \texttt{pau\_Latn} & Palauan & M \\
        \texttt{pbb\_Latn} & Páez & M \\
        \texttt{pbc\_Latn} & Patamona & M \\
        \texttt{pbi\_Latn} & Parkwa & M \\
        \texttt{pbs\_Latn} & Central Pame & L \\
        \texttt{pbt\_Arab} & Southern Pashto & L \\
        \texttt{pbu\_Arab} & Northern Pashto & L \\
        \texttt{pce\_Thai} & Ruching Palaung & M \\
        \texttt{pcm\_Latn} & Nigerian Pidgin & M \\
        \texttt{pex\_Latn} & Petats & L \\
        \texttt{pez\_Latn} & Eastern Penan & M \\
        \texttt{phl\_Arab} & Phalura & M \\
        \texttt{phr\_Arab} & Pahari-Potwari & M \\
        \texttt{pib\_Latn} & Yine & M \\
        \texttt{pil\_Latn} & Yom & M \\
        \texttt{pip\_Latn} & Pero & L \\
        \texttt{pir\_Latn} & Piratapuyo & M \\
        \texttt{pis\_Latn} & Pijin & M \\
        \texttt{piy\_Latn} & Piya-Kwonci & L \\
        \texttt{pjt\_Latn} & Pitjantjatjara & H \\
        \texttt{pkb\_Latn} & Pokomo & M \\
        \texttt{pko\_Latn} & Pökoot & L \\
        \texttt{plk\_Arab} & Kohistani Shina & M \\
        \texttt{pls\_Latn} & San Marcos Tlacoyalco Popoloca & M \\
        \texttt{plt\_Latn} & Plateau Malagasy & M \\
        \texttt{plw\_Latn} & Brooke's Point Palawano & M \\
        \texttt{pmf\_Latn} & Pamona & M \\
        \texttt{pmq\_Latn} & Northern Pame & L \\
        \texttt{pms\_Latn} & Piemontese & L \\
        \texttt{pmy\_Latn} & Papuan Malay & L \\
        \texttt{pnb\_Arab} & Western Panjabi & L \\
        \texttt{pne\_Latn} & Western Penan & L \\
        \texttt{pny\_Latn} & Pinyin & M \\
        \texttt{poc\_Latn} & Poqomam & L \\
        \texttt{poe\_Latn} & San Juan Atzingo Popoloca & L \\
        \texttt{poh\_Latn} & Poqomchi' & H \\
        \texttt{poi\_Latn} & Highland Popoluca & M \\
        \texttt{pol\_Latn} & Polish & H \\
        \texttt{por\_Latn} & Portuguese & H \\
        \texttt{pov\_Latn} & Upper Guinea Crioulo & L \\
        \texttt{pow\_Latn} & San Felipe Otlaltepec Popoloca & L \\
        \texttt{poy\_Latn} & Pogolo & M \\
        \texttt{ppk\_Latn} & Uma & H \\
        \texttt{pps\_Latn} & San Luís Temalacayuca Popoloca & M \\
        \texttt{prf\_Latn} & Paranan & M \\
        \texttt{prk\_Latn} & Parauk & M \\
        \texttt{prq\_Latn} & Ashéninka Perené & L \\
        \texttt{prt\_Thai} & Phai & H \\
        \texttt{pse\_Latn} & Central Malay & M \\
        \texttt{pss\_Latn} & Kaulong & H \\
        \texttt{pst\_Arab} & Central Pashto & H \\
        \texttt{ptu\_Latn} & Bambam & M \\
        \texttt{pua\_Latn} & Western Highland Purepecha & L \\
        \texttt{pui\_Latn} & Puinave & M \\
        \texttt{pus\_Arab} & Pushto & M \\
        \texttt{pwg\_Latn} & Gapapaiwa & H \\
        \texttt{pwn\_Latn} & Paiwan & M \\
        \texttt{pww\_Thai} & Pwo Northern Karen & M \\
        \texttt{pxm\_Latn} & Quetzaltepec Mixe & M \\
        \texttt{qub\_Latn} & Huallaga Huánuco Quechua & M \\
        \texttt{quc\_Latn} & K'iche' & H \\
        \texttt{quf\_Latn} & Lambayeque Quechua & M \\
        \texttt{qug\_Latn} & Chimborazo Highland Quichua & L \\
        \texttt{quh\_Latn} & South Bolivian Quechua & H \\
        \texttt{qul\_Latn} & North Bolivian Quechua & M \\
        \texttt{qum\_Latn} & Sipacapense & L \\
        \texttt{qup\_Latn} & Southern Pastaza Quechua & L \\
        \texttt{qur\_Latn} & Yanahuanca Pasco Quechua & L \\
        \texttt{qus\_Latn} & Santiago del Estero Quichua & L \\
        \texttt{quv\_Latn} & Sacapulteco & L \\
        \texttt{quw\_Latn} & Tena Lowland Quichua & M \\
        \texttt{qux\_Latn} & Yauyos Quechua & L \\
        \texttt{quy\_Latn} & Ayacucho Quechua & M \\
        \texttt{quz\_Latn} & Cusco Quechua & M \\
        \texttt{qva\_Latn} & Ambo-Pasco Quechua & L \\
        \texttt{qvc\_Latn} & Cajamarca Quechua & M \\
        \texttt{qve\_Latn} & Eastern Apurímac Quechua & M \\
        \texttt{qvh\_Latn} & Huamalíes-Dos de Mayo Huánuco Quechua & M \\
        \texttt{qvi\_Latn} & Imbabura Highland Quichua & L \\
        \texttt{qvj\_Latn} & Loja Highland Quichua & L \\
        \texttt{qvl\_Latn} & Cajatambo North Lima Quechua & L \\
        \texttt{qvm\_Latn} & Margos-Yarowilca-Lauricocha Quechua & M \\
        \texttt{qvn\_Latn} & North Junín Quechua & M \\
        \texttt{qvo\_Latn} & Napo Lowland Quechua & M \\
        \texttt{qvs\_Latn} & San Martín Quechua & M \\
        \texttt{qvw\_Latn} & Huaylla Wanca Quechua & M \\
        \texttt{qvz\_Latn} & Northern Pastaza Quichua & M \\
        \texttt{qwa\_Latn} & Corongo Ancash Quechua & L \\
        \texttt{qwh\_Latn} & Huaylas Ancash Quechua & M \\
        \texttt{qws\_Latn} & Sihuas Ancash Quechua & L \\
        \texttt{qxa\_Latn} & Chiquián Ancash Quechua & L \\
        \texttt{qxh\_Latn} & Panao Huánuco Quechua & M \\
        \texttt{qxl\_Latn} & Salasaca Highland Quichua & M \\
        \texttt{qxn\_Latn} & Northern Conchucos Ancash Quechua & M \\
        \texttt{qxo\_Latn} & Southern Conchucos Ancash Quechua & M \\
        \texttt{qxp\_Latn} & Puno Quechua & L \\
        \texttt{qxr\_Latn} & Cañar Highland Quichua & M \\
        \texttt{qxt\_Latn} & Santa Ana de Tusi Pasco Quechua & L \\
        \texttt{qxu\_Latn} & Arequipa-La Unión Quechua & L \\
        \texttt{qxw\_Latn} & Jauja Wanca Quechua & L \\
        \texttt{rag\_Latn} & Logooli & L \\
        \texttt{rah\_Beng} & Rabha & L \\
        \texttt{rai\_Latn} & Ramoaaina & M \\
        \texttt{rap\_Latn} & Rapanui & M \\
        \texttt{rav\_Deva} & Sampang & M \\
        \texttt{raw\_Latn} & Rawang & M \\
        \texttt{rej\_Latn} & Rejang & M \\
        \texttt{rel\_Latn} & Rendille & M \\
        \texttt{rgu\_Latn} & Ringgou & L \\
        \texttt{rhg\_Latn} & Rohingya & L \\
        \texttt{rif\_Arab} & Tarifit & M \\
        \texttt{rif\_Latn} & Tarifit & M \\
        \texttt{rim\_Latn} & Nyaturu & M \\
        \texttt{rjs\_Deva} & Rajbanshi & M \\
        \texttt{rkt\_Beng} & Rangpuri & M \\
        \texttt{rmc\_Cyrl} & Carpathian Romani & L \\
        \texttt{rmc\_Latn} & Carpathian Romani & M \\
        \texttt{rmo\_Latn} & Sinte Romani & M \\
        \texttt{rmy\_Cyrl} & Vlax Romani & L \\
        \texttt{rmy\_Latn} & Vlax Romani & M \\
        \texttt{rng\_Latn} & Ronga & M \\
        \texttt{rnl\_Latn} & Ranglong & M \\
        \texttt{rob\_Latn} & Tae' & L \\
        \texttt{rof\_Latn} & Rombo & M \\
        \texttt{roh\_Latn\_surs1244} & Romansh (Sursilvan) & L \\
        \texttt{rol\_Latn} & Romblomanon & M \\
        \texttt{ron\_Latn} & Romanian & H \\
        \texttt{roo\_Latn} & Rotokas & L \\
        \texttt{rop\_Latn} & Kriol & L \\
        \texttt{rro\_Latn} & Waima & M \\
        \texttt{rth\_Latn} & Ratahan & L \\
        \texttt{rub\_Latn} & Gungu & H \\
        \texttt{ruc\_Latn} & Ruuli & L \\
        \texttt{ruf\_Latn} & Luguru & M \\
        \texttt{rug\_Latn} & Roviana & M \\
        \texttt{run\_Latn} & Rundi & M \\
        \texttt{rus\_Cyrl} & Russian & H \\
        \texttt{rwm\_Latn} & Amba (Uganda) & L \\
        \bottomrule
    \end{tabular}
    \end{adjustbox}
}

{
    \tiny\renewcommand{\arraystretch}{.8}
    \begin{adjustbox}{max width=\textwidth}
    \begin{tabular}{llc}
        \toprule
        \textbf{Code} & \textbf{Name} & \textbf{Res}\\
        \midrule
        \texttt{rwr\_Deva} & Marwari (India) & L \\
        \texttt{sab\_Latn} & Buglere & M \\
        \texttt{sag\_Latn} & Sango & M \\
        \texttt{sah\_Cyrl} & Yakut & M \\
        \texttt{saj\_Latn} & Sahu & M \\
        \texttt{saq\_Latn} & Samburu & M \\
        \texttt{sas\_Latn} & Sasak & M \\
        \texttt{sau\_Latn} & Saleman & L \\
        \texttt{say\_Latn} & Saya & L \\
        \texttt{sba\_Latn} & Ngambay & M \\
        \texttt{sbd\_Latn} & Southern Samo & M \\
        \texttt{sbl\_Latn} & Botolan Sambal & M \\
        \texttt{sbn\_Arab} & Sindhi Bhil & L \\
        \texttt{sbp\_Latn} & Sangu (Tanzania) & H \\
        \texttt{sch\_Latn} & Sakachep & M \\
        \texttt{sck\_Deva} & Sadri & M \\
        \texttt{scl\_Arab} & Shina & L \\
        \texttt{scn\_Latn} & Sicilian & L \\
        \texttt{sco\_Latn} & Scots & L \\
        \texttt{sda\_Latn} & Toraja-Sa'dan & M \\
        \texttt{sdo\_Latn} & Bukar-Sadung Bidayuh & L \\
        \texttt{sea\_Latn} & Semai & L \\
        \texttt{seh\_Latn} & Sena & M \\
        \texttt{sei\_Latn} & Seri & L \\
        \texttt{ses\_Latn} & Koyraboro Senni Songhai & M \\
        \texttt{sey\_Latn} & Secoya & H \\
        \texttt{sgb\_Latn} & Mag-antsi Ayta & M \\
        \texttt{sgj\_Deva} & Surgujia & M \\
        \texttt{sgw\_Ethi} & Sebat Bet Gurage & M \\
        \texttt{shi\_Latn} & Tachelhit & M \\
        \texttt{shk\_Latn} & Shilluk & M \\
        \texttt{shn\_Mymr} & Shan & M \\
        \texttt{sho\_Latn} & Shanga & M \\
        \texttt{shp\_Latn} & Shipibo-Conibo & M \\
        \texttt{sid\_Latn} & Sidamo & M \\
        \texttt{sig\_Latn} & Paasaal & M \\
        \texttt{sil\_Latn} & Tumulung Sisaala & M \\
        \texttt{sin\_Sinh} & Sinhala & L \\
        \texttt{sip\_Tibt} & Sikkimese & L \\
        \texttt{siw\_Latn} & Siwai & L \\
        \texttt{sja\_Latn} & Epena & M \\
        \texttt{sjm\_Latn} & Mapun & M \\
        \texttt{sjp\_Deva} & Surjapuri & L \\
        \texttt{sjr\_Latn} & Siar-Lak & L \\
        \texttt{skg\_Latn} & Sakalava Malagasy & L \\
        \texttt{skr\_Arab} & Saraiki & L \\
        \texttt{sld\_Latn} & Sissala & M \\
        \texttt{slk\_Latn} & Slovak & H \\
        \texttt{slu\_Latn} & Selaru & L \\
        \texttt{slv\_Latn} & Slovenian & H \\
        \texttt{sml\_Latn} & Central Sama & M \\
        \texttt{smo\_Latn} & Samoan & M \\
        \texttt{sna\_Latn} & Shona & M \\
        \texttt{snc\_Latn} & Sinaugoro & L \\
        \texttt{snd\_Arab} & Sindhi & M \\
        \texttt{sne\_Latn} & Bau Bidayuh & M \\
        \texttt{snk\_Latn} & Soninke & L \\
        \texttt{snn\_Latn} & Siona & H \\
        \texttt{snp\_Latn} & Siane & M \\
        \texttt{snv\_Latn} & Sa'ban & L \\
        \texttt{snw\_Latn} & Selee & M \\
        \texttt{sol\_Latn} & Solos & L \\
        \texttt{som\_Latn} & Somali & H \\
        \texttt{soy\_Latn} & Miyobe & M \\
        \texttt{spa\_Latn} & Spanish & H \\
        \texttt{spp\_Latn} & Supyire Senoufo & M \\
        \texttt{sps\_Latn} & Saposa & L \\
        \texttt{spy\_Latn} & Sabaot & M \\
        \texttt{src\_Latn} & Logudorese Sardinian & L \\
        \texttt{srd\_Latn} & Sardinian & L \\
        \texttt{sri\_Latn} & Siriano & M \\
        \texttt{srm\_Latn} & Saramaccan & M \\
        \texttt{srn\_Latn} & Sranan Tongo & M \\
        \texttt{sro\_Latn} & Campidanese Sardinian & L \\
        \texttt{srp\_Cyrl} & Serbian & H \\
        \texttt{srr\_Latn} & Serer & L \\
        \texttt{srx\_Deva} & Sirmauri & M \\
        \texttt{ssi\_Arab} & Sansi & L \\
        \texttt{ste\_Latn} & Liana-Seti & L \\
        \texttt{stn\_Latn} & Owa & H \\
        \texttt{stp\_Latn} & Southeastern Tepehuan & M \\
        \texttt{sua\_Latn} & Sulka & L \\
        \texttt{suc\_Latn} & Western Subanon & M \\
        \texttt{suk\_Latn} & Sukuma & M \\
        \texttt{sun\_Latn} & Sundanese & H \\
        \texttt{sur\_Latn} & Mwaghavul & M \\
        \texttt{sus\_Latn} & Susu & M \\
        \texttt{suv\_Latn} & Puroik & L \\
        \texttt{suz\_Deva} & Sunwar & M \\
        \texttt{sva\_Geor} & Svan & M \\
        \texttt{swe\_Latn} & Swedish & H \\
        \texttt{swh\_Latn} & Swahili (individual language) & H \\
        \texttt{swv\_Deva} & Shekhawati & L \\
        \texttt{sxb\_Latn} & Suba & H \\
        \texttt{sxn\_Latn} & Sangir & M \\
        \texttt{sya\_Latn} & Siang & L \\
        \texttt{syl\_Latn} & Sylheti & L \\
        \texttt{sza\_Latn} & Semelai & L \\
        \texttt{szy\_Latn} & Sakizaya & M \\
        \texttt{tac\_Latn} & Lowland Tarahumara & M \\
        \texttt{taj\_Deva} & Eastern Tamang & M \\
        \texttt{tam\_Taml} & Tamil & H \\
        \texttt{tan\_Latn} & Tangale & L \\
        \texttt{tao\_Latn} & Yami & H \\
        \texttt{tap\_Latn} & Taabwa & M \\
        \texttt{taq\_Latn} & Tamasheq & M \\
        \texttt{tar\_Latn} & Central Tarahumara & L \\
        \texttt{tat\_Cyrl} & Tatar & M \\
        \texttt{tav\_Latn} & Tatuyo & H \\
        \texttt{tay\_Latn} & Atayal & L \\
        \texttt{tbc\_Latn} & Takia & M \\
        \texttt{tbf\_Latn} & Mandara & L \\
        \texttt{tbg\_Latn} & North Tairora & M \\
        \texttt{tbk\_Latn} & Calamian Tagbanwa & H \\
        \texttt{tbl\_Latn} & Tboli & H \\
        \texttt{tby\_Latn} & Tabaru & M \\
        \texttt{tbz\_Latn} & Ditammari & M \\
        \texttt{tca\_Latn} & Ticuna & M \\
        \texttt{tcc\_Latn} & Datooga & M \\
        \texttt{tcf\_Latn} & Malinaltepec Me'phaa & L \\
        \texttt{tcy\_Mlym} & Tulu & L \\
        \texttt{tcz\_Latn} & Thado Chin & L \\
        \texttt{tdj\_Latn} & Tajio & L \\
        \texttt{tdn\_Latn} & Tondano & L \\
        \texttt{tdx\_Latn} & Tandroy-Mahafaly Malagasy & L \\
        \texttt{ted\_Latn} & Tepo Krumen & M \\
        \texttt{tee\_Latn} & Huehuetla Tepehua & M \\
        \texttt{tel\_Telu} & Telugu & H \\
        \texttt{tem\_Latn} & Timne & M \\
        \texttt{teo\_Latn} & Teso & M \\
        \texttt{ter\_Latn} & Tereno & M \\
        \texttt{tew\_Latn} & Tewa (USA) & M \\
        \texttt{tex\_Latn} & Tennet & M \\
        \bottomrule
    \end{tabular}
    \quad
    \begin{tabular}{llc}
        \toprule
        \textbf{Code} & \textbf{Name} & \textbf{Res}\\
        \midrule
        \texttt{tfr\_Latn} & Teribe & M \\
        \texttt{tgc\_Latn} & Tigak & L \\
        \texttt{tgj\_Latn} & Tagin & H \\
        \texttt{tgk\_Cyrl} & Tajik & H \\
        \texttt{tgl\_Latn} & Tagalog & L \\
        \texttt{tgo\_Latn} & Sudest & M \\
        \texttt{tgp\_Latn} & Tangoa & M \\
        \texttt{tha\_Thai} & Thai & H \\
        \texttt{the\_Deva} & Chitwania Tharu & L \\
        \texttt{thk\_Latn} & Tharaka & M \\
        \texttt{thl\_Deva} & Dangaura Tharu & M \\
        \texttt{thq\_Deva} & Kochila Tharu & L \\
        \texttt{thr\_Deva} & Rana Tharu & L \\
        \texttt{thv\_Tfng} & Tahaggart Tamahaq & L \\
        \texttt{tig\_Ethi} & Tigre & L \\
        \texttt{tih\_Latn} & Timugon Murut & M \\
        \texttt{tik\_Latn} & Tikar & M \\
        \texttt{tio\_Latn} & Teop & L \\
        \texttt{tir\_Ethi} & Tigrinya & M \\
        \texttt{tkg\_Latn} & Tesaka Malagasy & M \\
        \texttt{tkr\_Latn} & Tsakhur & L \\
        \texttt{tkt\_Deva} & Kathoriya Tharu & L \\
        \texttt{tlb\_Latn} & Tobelo & H \\
        \texttt{tli\_Latn} & Tlingit & L \\
        \texttt{tlj\_Latn} & Talinga-Bwisi & M \\
        \texttt{tlp\_Latn} & Filomena Mata-Coahuitlán Totonac & L \\
        \texttt{tly\_Latn} & Talysh & M \\
        \texttt{tmc\_Latn} & Tumak & M \\
        \texttt{tmf\_Latn} & Toba-Maskoy & H \\
        \texttt{tna\_Latn} & Tacana & H \\
        \texttt{tng\_Latn} & Tobanga & M \\
        \texttt{tnk\_Latn} & Kwamera & M \\
        \texttt{tnn\_Latn} & North Tanna & M \\
        \texttt{tnp\_Latn} & Whitesands & L \\
        \texttt{tnr\_Latn} & Ménik & H \\
        \texttt{tnt\_Latn} & Tontemboan & H \\
        \texttt{tob\_Latn} & Toba & H \\
        \texttt{toc\_Latn} & Coyutla Totonac & M \\
        \texttt{toh\_Latn} & Gitonga & M \\
        \texttt{tok\_Latn} & Toki Pona & L \\
        \texttt{tom\_Latn} & Tombulu & M \\
        \texttt{top\_Latn} & Papantla Totonac & L \\
        \texttt{tos\_Latn} & Highland Totonac & M \\
        \texttt{tpi\_Latn} & Tok Pisin & H \\
        \texttt{tpl\_Latn} & Tlacoapa Me'phaa & L \\
        \texttt{tpm\_Latn} & Tampulma & M \\
        \texttt{tpp\_Latn} & Pisaflores Tepehua & M \\
        \texttt{tpt\_Latn} & Tlachichilco Tepehua & M \\
        \texttt{tpz\_Latn} & Tinputz & L \\
        \texttt{tqp\_Latn} & Tomoip & L \\
        \texttt{trc\_Latn} & Copala Triqui & M \\
        \texttt{tri\_Latn} & Trió & M \\
        \texttt{trn\_Latn} & Trinitario & M \\
        \texttt{trp\_Latn} & Kok Borok & L \\
        \texttt{trq\_Latn} & San Martín Itunyoso Triqui & L \\
        \texttt{trs\_Latn} & Chicahuaxtla Triqui & M \\
        \texttt{trv\_Latn} & Sediq & L \\
        \texttt{trw\_Arab} & Torwali & M \\
        \texttt{tsn\_Latn} & Tswana & L \\
        \texttt{tso\_Latn} & Tsonga & M \\
        \texttt{tsz\_Latn} & Purepecha & M \\
        \texttt{ttc\_Latn} & Tektiteko & H \\
        \texttt{tte\_Latn} & Bwanabwana & M \\
        \texttt{ttj\_Latn} & Tooro & M \\
        \texttt{ttq\_Tfng} & Tawallammat Tamajaq & M \\
        \texttt{ttr\_Latn} & Tera & L \\
        \texttt{ttu\_Latn} & Torau & L \\
        \texttt{tue\_Latn} & Tuyuca & M \\
        \texttt{tuf\_Latn} & Central Tunebo & H \\
        \texttt{tui\_Latn} & Tupuri & L \\
        \texttt{tuk\_Arab} & Turkmen & M \\
        \texttt{tuk\_Latn} & Turkmen & M \\
        \texttt{tul\_Latn} & Tula & L \\
        \texttt{tuo\_Latn} & Tucano & M \\
        \texttt{tuq\_Latn} & Tedaga & L \\
        \texttt{tur\_Latn} & Turkish & H \\
        \texttt{tuv\_Latn} & Turkana & L \\
        \texttt{tuy\_Latn} & Tugen & L \\
        \texttt{tvo\_Latn} & Tidore & L \\
        \texttt{tvu\_Latn} & Tunen & L \\
        \texttt{tvw\_Latn} & Sedoa & H \\
        \texttt{twb\_Latn} & Western Tawbuid & M \\
        \texttt{twe\_Latn} & Tewa (Indonesia) & L \\
        \texttt{twu\_Latn} & Termanu & M \\
        \texttt{txa\_Latn} & Tombonuo & M \\
        \texttt{txq\_Latn} & Tii & M \\
        \texttt{txs\_Latn} & Tonsea & L \\
        \texttt{txu\_Latn} & Kayapó & H \\
        \texttt{txy\_Latn} & Tanosy Malagasy & L \\
        \texttt{tye\_Latn} & Kyanga & M \\
        \texttt{tzh\_Latn} & Tzeltal & M \\
        \texttt{tzj\_Latn} & Tz'utujil & H \\
        \texttt{tzo\_Latn} & Tzotzil & M \\
        \texttt{ubl\_Latn} & Buhi'non Bikol & M \\
        \texttt{ubu\_Latn} & Umbu-Ungu & H \\
        \texttt{udl\_Latn} & Wuzlam & L \\
        \texttt{udm\_Cyrl} & Udmurt & M \\
        \texttt{udu\_Latn} & Uduk & M \\
        \texttt{uig\_Arab} & Uighur & H \\
        \texttt{uig\_Cyrl} & Uighur & M \\
        \texttt{uki\_Orya} & Kui (India) & L \\
        \texttt{ukr\_Cyrl} & Ukrainian & H \\
        \texttt{ukv\_Latn} & Kuku & L \\
        \texttt{umb\_Latn} & Umbundu & M \\
        \texttt{upv\_Latn} & Uripiv-Wala-Rano-Atchin & M \\
        \texttt{ura\_Latn} & Urarina & M \\
        \texttt{urb\_Latn} & Urubú-Kaapor & H \\
        \texttt{urd\_Arab} & Urdu & H \\
        \texttt{urd\_Deva} & Urdu & M \\
        \texttt{urd\_Latn} & Urdu & M \\
        \texttt{urh\_Latn} & Urhobo & L \\
        \texttt{urk\_Thai} & Urak Lawoi' & M \\
        \texttt{urt\_Latn} & Urat & H \\
        \texttt{ury\_Latn} & Orya & H \\
        \texttt{ush\_Arab} & Ushojo & L \\
        \texttt{usp\_Latn} & Uspanteco & M \\
        \texttt{uzb\_Cyrl} & Uzbek & H \\
        \texttt{uzb\_Latn} & Uzbek & H \\
        \texttt{uzn\_Latn} & Northern Uzbek & M \\
        \texttt{vag\_Latn} & Vagla & M \\
        \texttt{vah\_Deva} & Varhadi-Nagpuri & L \\
        \texttt{vai\_Latn} & Vai & L \\
        \texttt{var\_Latn} & Huarijio & L \\
        \texttt{ver\_Latn} & Mom Jango & L \\
        \texttt{vid\_Latn} & Vidunda & M \\
        \texttt{vie\_Latn} & Vietnamese & H \\
        \texttt{vif\_Latn} & Vili & M \\
        \texttt{vmc\_Latn} & Juxtlahuaca Mixtec & L \\
        \texttt{vmj\_Latn} & Ixtayutla Mixtec & L \\
        \texttt{vmm\_Latn} & Mitlatongo Mixtec & L \\
        \texttt{vmp\_Latn} & Soyaltepec Mazatec & L \\
        \texttt{vmw\_Latn} & Makhuwa & M \\
        \texttt{vmy\_Latn} & Ayautla Mazatec & M \\
        \bottomrule
    \end{tabular}
    \quad
    \begin{tabular}{llc}
        \toprule
        \textbf{Code} & \textbf{Name} & \textbf{Res}\\
        \midrule
        \texttt{vmz\_Latn} & Mazatlán Mazatec & L \\
        \texttt{vro\_Latn} & Võro & L \\
        \texttt{vun\_Latn} & Vunjo & M \\
        \texttt{vut\_Latn} & Vute & M \\
        \texttt{wal\_Ethi} & Wolaytta & M \\
        \texttt{wal\_Latn} & Wolaytta & M \\
        \texttt{wap\_Latn} & Wapishana & H \\
        \texttt{war\_Latn} & Waray (Philippines) & M \\
        \texttt{waw\_Latn} & Waiwai & M \\
        \texttt{way\_Latn} & Wayana & M \\
        \texttt{wba\_Latn} & Warao & M \\
        \texttt{wbl\_Latn} & Wakhi & L \\
        \texttt{wbr\_Deva} & Wagdi & L \\
        \texttt{wci\_Latn} & Waci Gbe & L \\
        \texttt{weo\_Latn} & Wemale & L \\
        \texttt{wes\_Latn} & Cameroon Pidgin & L \\
        \texttt{wja\_Latn} & Waja & L \\
        \texttt{wji\_Latn} & Warji & L \\
        \texttt{wlo\_Latn} & Wolio & L \\
        \texttt{wlx\_Latn} & Wali (Ghana) & M \\
        \texttt{wmw\_Latn} & Mwani & M \\
        \texttt{wob\_Latn} & Wè Northern & M \\
        \texttt{wof\_Latn} & Gambian Wolof & L \\
        \texttt{wol\_Latn} & Wolof & M \\
        \texttt{wsg\_Telu} & Adilabad Gondi & M \\
        \texttt{wwa\_Latn} & Waama & M \\
        \texttt{xal\_Cyrl} & Kalmyk & M \\
        \texttt{xdy\_Latn} & Malayic Dayak & L \\
        \texttt{xed\_Latn} & Hdi & M \\
        \texttt{xer\_Latn} & Xerénte & L \\
        \texttt{xhe\_Arab} & Khetrani & L \\
        \texttt{xho\_Latn} & Xhosa & M \\
        \texttt{xka\_Arab} & Kalkoti & L \\
        \texttt{xkl\_Latn} & Mainstream Kenyah & L \\
        \texttt{xmf\_Geor} & Mingrelian & L \\
        \texttt{xmm\_Latn} & Manado Malay & H \\
        \texttt{xmv\_Latn} & Antankarana Malagasy & M \\
        \texttt{xnj\_Latn} & Ngoni (Tanzania) & M \\
        \texttt{xnr\_Deva} & Kangri & M \\
        \texttt{xog\_Latn} & Soga & M \\
        \texttt{xon\_Latn} & Konkomba & M \\
        \texttt{xpe\_Latn} & Liberia Kpelle & L \\
        \texttt{xrb\_Latn} & Eastern Karaboro & M \\
        \texttt{xsb\_Latn} & Sambal & M \\
        \texttt{xsm\_Latn} & Kasem & M \\
        \texttt{xsr\_Deva} & Sherpa & M \\
        \texttt{xsu\_Latn} & Sanumá & M \\
        \texttt{xta\_Latn} & Alcozauca Mixtec & L \\
        \texttt{xtd\_Latn} & Diuxi-Tilantongo Mixtec & H \\
        \texttt{xte\_Latn} & Ketengban & H \\
        \texttt{xti\_Latn} & Sinicahua Mixtec & L \\
        \texttt{xtm\_Latn} & Magdalena Peñasco Mixtec & H \\
        \texttt{xtn\_Latn} & Northern Tlaxiaco Mixtec & M \\
        \texttt{xtu\_Latn} & Cuyamecalco Mixtec & L \\
        \texttt{xua\_Taml} & Alu Kurumba & L \\
        \texttt{xuo\_Latn} & Kuo & M \\
        \texttt{yaa\_Latn} & Yaminahua & M \\
        \texttt{yad\_Latn} & Yagua & M \\
        \texttt{yal\_Latn} & Yalunka & M \\
        \texttt{yam\_Latn} & Yamba & M \\
        \texttt{yao\_Latn} & Yao & M \\
        \texttt{yaq\_Latn} & Yaqui & L \\
        \texttt{yas\_Latn} & Nugunu (Cameroon) & M \\
        \texttt{yat\_Latn} & Yambeta & M \\
        \texttt{yav\_Latn} & Yangben & L \\
        \texttt{yay\_Latn} & Agwagwune & L \\
        \texttt{yaz\_Latn} & Lokaa & H \\
        \texttt{yba\_Latn} & Yala & M \\
        \texttt{ybb\_Latn} & Yemba & M \\
        \texttt{ycl\_Latn} & Lolopo & H \\
        \texttt{ycn\_Latn} & Yucuna & M \\
        \texttt{ydd\_Hebr} & Eastern Yiddish & M \\
        \texttt{ydg\_Arab} & Yidgha & L \\
        \texttt{yea\_Mlym} & Ravula & M \\
        \texttt{yer\_Latn} & Tarok & L \\
        \texttt{yes\_Latn} & Nyankpa & L \\
        \texttt{yka\_Latn} & Yakan & M \\
        \texttt{yli\_Latn} & Angguruk Yali & M \\
        \texttt{yor\_Latn} & Yoruba & H \\
        \texttt{yre\_Latn} & Yaouré & M \\
        \texttt{yua\_Latn} & Yucateco & M \\
        \texttt{yue\_Hans} & Yue Chinese & H \\
        \texttt{yue\_Hant} & Yue Chinese & M \\
        \texttt{yuz\_Latn} & Yuracare & M \\
        \texttt{yva\_Latn} & Yawa & M \\
        \texttt{zaa\_Latn} & Sierra de Juárez Zapotec & M \\
        \texttt{zab\_Latn} & Western Tlacolula Valley Zapotec & M \\
        \texttt{zac\_Latn} & Ocotlán Zapotec & L \\
        \texttt{zad\_Latn} & Cajonos Zapotec & M \\
        \texttt{zae\_Latn} & Yareni Zapotec & M \\
        \texttt{zai\_Latn} & Isthmus Zapotec & M \\
        \texttt{zam\_Latn} & Miahuatlán Zapotec & L \\
        \texttt{zao\_Latn} & Ozolotepec Zapotec & M \\
        \texttt{zaq\_Latn} & Aloápam Zapotec & H \\
        \texttt{zar\_Latn} & Rincón Zapotec & M \\
        \texttt{zas\_Latn} & Santo Domingo Albarradas Zapotec & M \\
        \texttt{zav\_Latn} & Yatzachi Zapotec & L \\
        \texttt{zaw\_Latn} & Mitla Zapotec & M \\
        \texttt{zca\_Latn} & Coatecas Altas Zapotec & M \\
        \texttt{zga\_Latn} & Kinga & H \\
        \texttt{zim\_Latn} & Mesme & M \\
        \texttt{ziw\_Latn} & Zigula & M \\
        \texttt{zmz\_Latn} & Mbandja & M \\
        \texttt{zne\_Latn} & Zande (individual language) & M \\
        \texttt{zoc\_Latn} & Copainalá Zoque & L \\
        \texttt{zoh\_Latn} & Chimalapa Zoque & L \\
        \texttt{zor\_Latn} & Rayón Zoque & L \\
        \texttt{zos\_Latn} & Francisco León Zoque & M \\
        \texttt{zpc\_Latn} & Choapan Zapotec & H \\
        \texttt{zpg\_Latn} & Guevea De Humboldt Zapotec & M \\
        \texttt{zpi\_Latn} & Santa María Quiegolani Zapotec & M \\
        \texttt{zpl\_Latn} & Lachixío Zapotec & M \\
        \texttt{zpm\_Latn} & Mixtepec Zapotec & M \\
        \texttt{zpo\_Latn} & Amatlán Zapotec & M \\
        \texttt{zpt\_Latn} & San Vicente Coatlán Zapotec & M \\
        \texttt{zpu\_Latn} & Yalálag Zapotec & M \\
        \texttt{zpv\_Latn} & Chichicapan Zapotec & L \\
        \texttt{zpy\_Latn} & Mazaltepec Zapotec & L \\
        \texttt{zpz\_Latn} & Texmelucan Zapotec & M \\
        \texttt{zsm\_Latn} & Standard Malay & H \\
        \texttt{ztg\_Latn} & Xanaguía Zapotec & L \\
        \texttt{ztn\_Latn} & Santa Catarina Albarradas Zapotec & L \\
        \texttt{ztp\_Latn} & Loxicha Zapotec & L \\
        \texttt{ztq\_Latn} & Quioquitani-Quierí Zapotec & M \\
        \texttt{zts\_Latn} & Tilquiapan Zapotec & L \\
        \texttt{ztu\_Latn} & Güilá Zapotec & L \\
        \texttt{zty\_Latn} & Yatee Zapotec & M \\
        \texttt{zul\_Latn} & Zulu & H \\
        \texttt{zyb\_Latn} & Yongbei Zhuang & M \\
        \texttt{zyp\_Latn} & Zyphe Chin & M \\
        \texttt{zza\_Latn} & Zaza & L \\
        \bottomrule
    \end{tabular}
    \end{adjustbox}
}
\captionof{table}{Full list of languages supported by \OmniASR, including language code, English name, and resource level (Low, Medium, High).}\label{tab:full_lang_list}

\section{WER Filtering}
\label{sec:appendix_wer_filtering}

WER-thresholds were used to filter out samples likely to be of low quality from the \vendorcorpus{} ASR dataset. Values ranged from 150 to 250 WER. These were determined qualitatively and selected to filter out samples with obviously misaligned audio/text. For example:

\hfill

\begin{verbatim}
Reference:
okoro ekwup mmotima nson wo mawanne ochike machip akpan pimoruku bebogye
Hypothesis:
okoro ekwu otok kpena kpe fu bok obo mo tim so woma wane mo chike ma achit
akpe pa mo orugo be boya bep be bae bake bonga akpe pe nok boya

Reference:
en sa w konn sa k pase
Hypothesis:
en fò w konn sa k pase n ap tou benefisye yon staj men m byen kwè so kò
kòman kote sa ye lankò menm chak ki bay bon moun yo wi me nm ja ou ka

Reference:
enh se fèt dè mè se fèt ou ankò
Hypothesis:
elepicit m konnen lepichit m konnen wi m konnen demis li rele en skisoee
bon tetout fason pann fèt aa o byen pete ye e fèe fèt b  èmè pis fèt ou ankò
\end{verbatim}

\hfill

In the above examples, it is clear in listening to the audio that the hypotheses generated by our model are more accurate than the reference texts, so we filtered such examples out.

\section{Prompts and Guidelines for Commissioned Data Collection}
This section contains the recording prompts and transcription guidelines for our commissioned data collection.
\subsection{Recording guidelines}\label{appx:guideline:record}
\begin{itemize}
    \item Please record in a quiet environment.
    \item During the recording, please refrain from: 
    \begin{itemize}
        \item touching the microphone,
        \item blowing into the microphone,
        \item moving things around that are close to the recording device.
    \end{itemize}
    \item Please refrain from clearing your throat, coughing, sneezing, or making any loud sounds during the recording.
    \item Please refrain from eating or drinking during the recording.
    \item Please speak in a natural, normal voice.
    \item Please speak at a normal pace and not too quickly or too slowly.
    \item If you encounter names and words that are in a different language (for example, an English name when you are speaking Swahili), please do your best to pronounce the name as you normally would in the target language.
    \item Please refrain from sharing any personally identifiable information in the recordings, whether it pertains to you or others. Personally identifiable information includes:
    \begin{itemize}
        \item Full name
        \item Phone number
        \item Home address
        \item Email address or other account identifiers, such as social media handles
        \item Passport number
        \item Social Security Number or analogous identification numbers
        \item Health information 
        \item Sexual orientation
        \item Political affiliation
        \item Any other analogous information
    \end{itemize}
\end{itemize}

\subsection{Transcription guidelines}\label{appx:guideline:transcript}
Your job is to transcribe exactly what was said in the recording, including a representation of all the disfluencies and noises it contains.
\begin{itemize}
    \item If the recording contains grammatical mistakes, these should not be corrected in the transcription.
    \item The only characters allowed in the transcription are letters of the given language, punctuation and the set of special tags specified below.
    \item (Updated) Wherever possible and if this is applicable to your language, please use punctuation in transcripts as you would normally do in your written language. Please also capitalize the beginnings of new sentences if applicable.
\end{itemize}

\paragraph{Numbers and acronyms.}
\begin{itemize}
    \item Numbers should be spelled out in words. They should not be written in the numeral system.
    \begin{itemize}
        \item \textcolor{red}{\textbf{Incorrect}}: \textit{I walked exactly 2017 steps.}
        \item \textcolor{ForestGreen}{\textbf{Correct}}: \textit{I walked exactly two thousand seventeen steps.}
    \end{itemize}
    \item Acronyms should be written as they are normally written in the language, following standard capitalization rules. They should not be transcribed phonetically.
    \begin{itemize}
        \item \textcolor{red}{\textbf{Incorrect}}: \textit{They were arrested by the eff bee eye last Thursday.}
        \item \textcolor{ForestGreen}{\textbf{Correct}}: \textit{They were arrested by the FBI last Thursday.}
    \end{itemize}
\end{itemize}

\paragraph{Punctuation and symbols}
\begin{itemize}
    \item Use the punctuation that is appropriate for writing in the given language.
    \item Symbols for currencies, percentages, etc. should be avoided, and should instead be spelled out.
    \begin{itemize}
        \item  \textcolor{red}{\textbf{Incorrect}}: \textit{This bag cost me only \$10!}
        \item \textcolor{ForestGreen}{\textbf{Correct}}: \textit{This bag cost me only ten dollars!}
    \end{itemize}
\end{itemize}

\paragraph{Special tags}

The following special tags should be used to mark disfluencies, fillers, and other types of non-verbal content.
\begin{table}[!ht]
    \centering
    \begin{tabular}{p{3cm}p{12cm}}
    \toprule
    \textbf{Tag} & \textbf{Meaning} \\ \midrule
    \texttt{<laugh>} & The sound of laughter. \\ 
\midrule
\texttt{<hesitation>} & A hesitation sound, often used by speakers while thinking of the next thing to say. In English, some common hesitation sounds are “err”, “um”, “huh”, etc. \\ 
\midrule
\texttt{<unintelligible>} & A word or sequence of words that cannot be understood.\\
\midrule
\texttt{<noise>} & Any other type of noise, such as the speaker coughing or clearing their throat, a car honking, the sound of something hitting the microphone, a phone buzzing, etc. \\
    \bottomrule
    \end{tabular}
    \caption{Special tags used for transcription}
    \label{tab:transcript_tags}
\end{table}
\begin{itemize}
    \item Tags should be inserted in the transcription at the appropriate location, and should be separated from the other content by spaces; for example:
    \begin{itemize}
        \item \textit{And then I \texttt{<noise>} went on holiday.}
        \item \textit{Well, \texttt{<noise>} \texttt{<laugh>} it wasn’t exactly a holiday \texttt{<laugh>}}
    \end{itemize}

\item When we speak, we often insert hesitations while thinking of the next idea we want to say. Some common hesitations in English are “err”, “um” and “uh”. Since these hesitations can vary significantly in the exact sounds and length used, and often there are no clear rules on how they should be written, for this project they should all be represented using the tag \texttt{<hesitation>}. Only this tag should be used. You should not attempt to transcribe hesitations using letters, such as “err”.
\end{itemize}

\paragraph{Word segments, false starts and repeated words.} \begin{itemize}
    \item Spontaneous speech naturally contains false starts where only a fragment of a full word is produced. For these instances, please transcribe to the best of your ability the word fragment and attach a hyphen at the end of the word (-) to indicate the word is a false start. 
    \begin{itemize}
        \item \textit{His name is \underline{Jo- Jona-} Jonathan.}
    \end{itemize}
    \item Sometimes speakers will repeat a word or word fragment multiple times. This should be transcribed too.
    \begin{itemize}
        \item \textit{And then I went to \underline{the the the bed-} the bedroom.}
    \end{itemize}

\end{itemize}


\paragraph{Grammatical mistakes and colloquialisms.}
\begin{itemize}
    \item Spontaneous speech will naturally contain grammatical mistakes. These should not be corrected when transcribing. The transcription should reflect the spoken content exactly.
    \item Speakers may use colloqualisms (such as, in English, “gonna”, “cuz”, etc.) which may not be considered formally correct. These should be transcribed as they are, and not changed to their more formal equivalents.
\end{itemize}

\section{Quality Assurance (QA) Guidelines}\label{appx:guideline:qa}
In this appendix, we detail the guidance provided to perform quality assurance (QA).
\subsection{Speech recording error taxonomy}
\Cref{tab:speech_qa_errors} shows the definitions used for each of the error categories. More broadly, QA technicians were asked to pay particular attention to the following speech recording issues:
\begin{itemize}
    \item General audio quality issues (e.g., volume is too low, speech is inaudible, there is constant background noise or heavy static, files seem systematically cut off before the end)
\item Ad hoc noises (e.g., rooster crowing, mechanical noises, bells or phones ringing, very long silences or pauses)
\item Other human voices (e.g., people talking in the background in the same language, or more problematic, in a different language)
\item The speaker responds to the prompts in a pivot language, not in the expected language (prompts were translated into a number of high-resource pivot languages and it can happen that the speaker will respond to the prompts in the same language as the prompts instead of responding in their native language)

\end{itemize}

\begin{table}[!ht]
    \centering
    \begin{tabular}{p{3cm}p{6cm}p{6cm}}
    \toprule
     \textbf{Category}  & \textbf{Critical example} & \textbf{Minor example} \\
     \midrule
     \small{Human vocal noise} & \small{Second voice in the background} & \small{N/A (This error is always critical)} \\
     & \small{Singing in the background} &  \\
    \midrule
    \small{Cutoff} & \small{Speech is cut off at either end of the recording} &  \small{N/A (This error is always critical)}\\
     \midrule
     \small{Background noise} & \small{Rooster crowing} & \small{Occasional mild coughing} \\
     & \small{Street noise, car honking} & \small{Occasional mild coughing} \\
     & \small{Bird chirping} & \small{Mild breathing sound}\\
     & \small{Strong wind} & \\
     \midrule
\small{Audio Glitches} & \small{Serious glitches that break up speech} & \small{Mild glitch happens in between speech} \\
\midrule
\small{Static noise} & \small{Strong static noise that affects intelligibility} & \small{Mild static noise that does not affect speech} \\
\midrule
\small{Low volume} & \small{Cannot hear the speech clearly in the max volume setting} & \small{Lower than normal but still audible at max volume} \\
\midrule
\small{Inconsistent volume} & \small{Volume changing drastically} & \small{Occasional soft voice} \\
\midrule
\small{Muffled voice} & \small{Muffled voice sounds like talking behind a curtain} & \small{Audio is not crisp but does not affect intelligibility} \\
\midrule
\small{Echo} & \small{Strong echo like speaking in a cave or tunnel such that it compromises the intelligibility of words} & \small{Mild echo in non-studio environment} \\
\midrule
\small{Microphone Noise} & \small{Any hissing, plosive, popping noise that breaks the speech} & \small{Mild pop noise when turning on/off the recorder} \\
\midrule
\small{Pause / Silence} & \small{Long pauses} & \small{Short pauses when speaker is thinking} \\
 & \small{- If at the start or end of speech and above 2s} & \\
 & \small{- If at the middle of speech and above 5s} & \\
 & \small{- If more than $\frac{1}{3}$ of the audio is made up of leading/trailing silence or intra-sentential silence (excluding normal pauses between words)} & \\
\midrule
\small{Unnatural speech} & \small{Consistent stutter or mumbling} & \small{Occasional repeated words and syllable} \\
 & \small{Extremely not fluent, words uttered individually} & \\
 & \small{Whisper} & \\
 & \small{Feels like someone reading / monotonous speech} & \\
     \bottomrule
    \end{tabular}
    \caption{Description of all error categories used for speech recording in-depth quality assurance.}
    \label{tab:speech_qa_errors}
\end{table}


\subsection{Transcript error taxonomy}

\Cref{tab:transcript_qa_errors} shows the definitions used for each of the error categories. More broadly, QA technicians were asked to pay particular attention to the following transcript issues:
\begin{itemize}
\item General transcription issues (e.g., the transcript does not match the audio file at all, the transcript is in an unexpected writing system, the transcript is in the International Phonetic Alphabet, the transcript is missing words, the transcript is much shorter or longer than it should be)
\item Transcription issues that are specific to a language (e.g., a few non-Unicode-compliant characters have been used)
\item Issues related to the use of event-marking tags (a specific tag set has been defined by the project team; \Cref{tab:transcript_tags})
\end{itemize}
\begin{table}[!ht]
    \centering
    \begin{tabular}{p{3.5cm}p{6cm}p{5.5cm}}
    \toprule
     \textbf{Category}  & \textbf{Critical example} & \textbf{Minor example} \\
     \midrule
    \small{Mismatch} & \small{Transcript file does not match the audio at all (either in content or in length) } & \small{N/A (mismatch is critical)} \\
\midrule
\small{Wrong writing system} & \small{The transcript does not use the expected writing system}  & \small{N/A (writing system is critical)} \\
 & \small{The transcript is in IPA or other phonetically-based system} & \\
 & \small{Different writing standard, inconsistency in the spelling (the same word spelled in different ways)} & \\
\midrule
\small{Wrong tags} & \small{The transcript includes made-up tags} & \small{N/A (all mistaggings are critical)} \\
 & \small{Tags are not used adequately (e.g., \texttt{<noise>} instead of \texttt{<hesitation>})} & \\
\midrule
\small{Numbers} & \small{The presence of numbers written in digits} & \small{(N/A writing digits is critical)} \\
\midrule
\small{Incomplete} &  \small{The transcript is abridged rather than verbatim} & \small{The transcript seems to sometimes be missing a word or two} \\
 & \small{The transcript consistently misses words} \\
\midrule
\small{Inconsistent tagging} & \small{The tag set being used is compliant but the transcriber consistently switches between tags for the same audio events} & \small{A few tags show inconsistency, especially for borderline audio events} \\

     \bottomrule
    \end{tabular}
    \caption{Description of all error categories used for transcript in-depth quality assurance.}
    \label{tab:transcript_qa_errors}
\end{table}